\documentclass[11pt,a4paper]{article}
\usepackage[utf8]{inputenc}
\usepackage{lmodern}
% Latin Modern: a scalable Type 1 font. With T1 and the default
% Computer Modern, TeX asks for the EC fonts, and when the Type 1 EC
% ones (cm-super) are missing it falls back to bitmaps -- and
% microtype's font expansion cannot work with a bitmap font, so the
% build fails on a TeX Live without cm-super.
\usepackage[T1]{fontenc}
\usepackage[english]{babel}
\usepackage[margin=2.3cm]{geometry}
\usepackage{xcolor}
\usepackage{booktabs}
\usepackage{longtable}
\usepackage{array}
\usepackage{amsmath,amssymb}
\usepackage{listings}
\usepackage{pdflscape}
\setlength{\emergencystretch}{6em}
\usepackage[hidelinks]{hyperref}
\newcommand{\code}[1]{\texttt{#1}}
\newcommand{\codex}[1]{\texttt{#1}}
% Generated by tools/gen_concepts_values.py -- do not edit.
% Values come from mal-core/config/model_constants.yaml via the
% validated loader; a key that no longer exists makes \mskey fail
% the build with an error, in any document that references it.

\makeatletter
\newcommand{\mskey}[1]{%
  \expandafter\ifx\csname ms@k@#1\endcsname\relax
    \errmessage{Concept document: unknown model constant key `#1'}%
  \else
    \csname ms@k@#1\endcsname%
  \fi}
\makeatother

\expandafter\def\csname ms@k@adults.daily_mort_basal\endcsname{0.87}
\expandafter\def\csname ms@k@adults.mort_cap\endcsname{0.93}
\expandafter\def\csname ms@k@adults.mort_floor\endcsname{0.6}
\expandafter\def\csname ms@k@adults.opt_c\endcsname{25.0}
\expandafter\def\csname ms@k@adults.sigma_c\endcsname{15.0}
\expandafter\def\csname ms@k@aquatic_stages.egg_daily_mort\endcsname{0.08}
\expandafter\def\csname ms@k@aquatic_stages.larva_bh_alpha\endcsname{0.05}
\expandafter\def\csname ms@k@aquatic_stages.larva_bh_s0\endcsname{0.95}
\expandafter\def\csname ms@k@aquatic_stages.last_larval_instar\endcsname{4}
\expandafter\def\csname ms@k@aquatic_stages.pupa_daily_mort\endcsname{0.05}
\expandafter\def\csname ms@k@aquatic_stages.salinity_missing_psu\endcsname{0.0}
\expandafter\def\csname ms@k@aquatic_stress.desicc_base_daily\endcsname{0.1}
\expandafter\def\csname ms@k@aquatic_stress.desicc_grace_days\endcsname{5.0}
\expandafter\def\csname ms@k@aquatic_stress.desicc_ramp_decay\endcsname{0.3}
\expandafter\def\csname ms@k@aquatic_stress.vulnerable_instar_max\endcsname{2}
\expandafter\def\csname ms@k@aquatic_stress.wash_fraction_max\endcsname{0.6}
\expandafter\def\csname ms@k@aquatic_stress.wash_threshold_mm\endcsname{40.0}
\expandafter\def\csname ms@k@capacity.persistent_availability\endcsname{1.0}
\expandafter\def\csname ms@k@contracts.cogs_state\endcsname{"4.0"}
\expandafter\def\csname ms@k@contracts.env_nc\endcsname{"3.0"}
\expandafter\def\csname ms@k@contracts.habitat\endcsname{"4.0"}
\expandafter\def\csname ms@k@dispersal.flight_hour\endcsname{20}
\expandafter\def\csname ms@k@dispersal.night_end_hour\endcsname{6}
\expandafter\def\csname ms@k@dispersal.night_start_hour\endcsname{18}
\expandafter\def\csname ms@k@dispersal.snap_radius_m\endcsname{500.0}
\expandafter\def\csname ms@k@dispersal.wind_flight_hours\endcsname{4.0}
\expandafter\def\csname ms@k@dispersal.wind_migration_prob\endcsname{0.05}
\expandafter\def\csname ms@k@dispersal.wind_self_speed_ms\endcsname{1.0}
\expandafter\def\csname ms@k@dispersal.wind_survival\endcsname{0.85}
\expandafter\def\csname ms@k@eip.base_c\endcsname{16.0}
\expandafter\def\csname ms@k@eip.threshold_gd\endcsname{110.0}
\expandafter\def\csname ms@k@geodesy.earth_radius_m\endcsname{6371008.8}
\expandafter\def\csname ms@k@geodesy.metres_per_deg_lat\endcsname{111195.08023353292}
\expandafter\def\csname ms@k@gonotrophic.cycle_duration_days\endcsname{2.0}
\expandafter\def\csname ms@k@gonotrophic.egg_batch_max\endcsname{170}
\expandafter\def\csname ms@k@gonotrophic.egg_batch_mean\endcsname{120}
\expandafter\def\csname ms@k@gonotrophic.egg_batch_min\endcsname{30}
\expandafter\def\csname ms@k@gonotrophic.egg_batch_n\endcsname{240}
\expandafter\def\csname ms@k@gonotrophic.feeding_success_rate\endcsname{0.825}
\expandafter\def\csname ms@k@gonotrophic.first_cycle_days\endcsname{4.0}
\expandafter\def\csname ms@k@gonotrophic.resting_duration_days\endcsname{0.5}
\expandafter\def\csname ms@k@habitat.activation_min_engine\endcsname{0.01}
\expandafter\def\csname ms@k@habitat.candidate_min_engine\endcsname{0.01}
\expandafter\def\csname ms@k@habitat_types.dynamics.crop_pluvial\endcsname{"pluvial"}
\expandafter\def\csname ms@k@habitat_types.dynamics.open_pluvial\endcsname{"pluvial"}
\expandafter\def\csname ms@k@habitat_types.dynamics.surface_persistent\endcsname{"persistent"}
\expandafter\def\csname ms@k@habitat_types.dynamics.urban_persistent\endcsname{"persistent"}
\expandafter\def\csname ms@k@habitat_types.dynamics.wetland_persistent\endcsname{"persistent"}
\expandafter\def\csname ms@k@habitat_types.dynamics.wetland_pluvial\endcsname{"pluvial"}
\expandafter\def\csname ms@k@habitat_types.types\endcsname{"surface\_persistent, wetland\_persistent, urban\_persistent, wetland\_pluvial, open\_pluvial, crop\_pluvial"}
\expandafter\def\csname ms@k@host_seeking.indoor_attraction_add\endcsname{0.72}
\expandafter\def\csname ms@k@host_seeking.indoor_fraction_rural\endcsname{0.3}
\expandafter\def\csname ms@k@host_seeking.indoor_fraction_urban\endcsname{0.72}
\expandafter\def\csname ms@k@host_seeking.starvation_days\endcsname{7}
\expandafter\def\csname ms@k@host_seeking.urbanicity_add\endcsname{0.2}
\expandafter\def\csname ms@k@ingest.candidate_min_ingest\endcsname{1e-06}
\expandafter\def\csname ms@k@ingest.cell_size_hint_m\endcsname{1000.0}
\expandafter\def\csname ms@k@ingest.coastline_buffer_m\endcsname{2000.0}
\expandafter\def\csname ms@k@ingest.coastline_resolution\endcsname{"f"}
\expandafter\def\csname ms@k@ingest.fill_frac_valid_max\endcsname{0.5}
\expandafter\def\csname ms@k@ingest.grid_degrees_step_max\endcsname{1}
\expandafter\def\csname ms@k@ingest.land_mask_threshold\endcsname{0.5}
\expandafter\def\csname ms@k@ingest.min_dilation_radius_cells\endcsname{1}
\expandafter\def\csname ms@k@ingest.nc_compression_level\endcsname{4}
\expandafter\def\csname ms@k@ingest.nodata_sentinel\endcsname{-9999.0}
\expandafter\def\csname ms@k@ingest.nodata_sentinel_int\endcsname{-9999}
\expandafter\def\csname ms@k@ingest.ocean_mask_threshold\endcsname{0.5}
\expandafter\def\csname ms@k@ingest.residential_from_indoor_frac\endcsname{0.8}
\expandafter\def\csname ms@k@ingest.sea_water_floor\endcsname{0.05}
\expandafter\def\csname ms@k@ingest.valid_count_epsilon\endcsname{1e-09}
\expandafter\def\csname ms@k@integration.directed_flight_budget_h\endcsname{0.4856}
\expandafter\def\csname ms@k@integration.dt_olfactory_s\endcsname{0.25}
\expandafter\def\csname ms@k@integration.dt_terminal_s\endcsname{0.04}
\expandafter\def\csname ms@k@integration.landing_speed_hold_s\endcsname{3.0}
\expandafter\def\csname ms@k@integration.landing_speed_mm_s\endcsname{50.0}
\expandafter\def\csname ms@k@integration.phase_duration_hours\endcsname{"12.0, 4.0, 6.0, 2.0"}
\expandafter\def\csname ms@k@integration.plume_memory_s\endcsname{0.7}
\expandafter\def\csname ms@k@mating.female_attendance_rate\endcsname{0.1033}
\expandafter\def\csname ms@k@mating.local_swarm_male_offer_factor\endcsname{0.0}
\expandafter\def\csname ms@k@mating.male_females_per_night_max\endcsname{3}
\expandafter\def\csname ms@k@mating.male_move_capacity_m\endcsname{50.0}
\expandafter\def\csname ms@k@mating.male_swarm_reference_size\endcsname{100.0}
\expandafter\def\csname ms@k@mating.p_mate_male_fraction\endcsname{0.072}
\expandafter\def\csname ms@k@mating.p_mate_pairs_per_individual\endcsname{0.167}
\expandafter\def\csname ms@k@mating.p_mate_size_exponent\endcsname{0.0}
\expandafter\def\csname ms@k@mating.rain_suppresses_swarming\endcsname{true}
\expandafter\def\csname ms@k@mating.swarm_duration_min\endcsname{28.3}
\expandafter\def\csname ms@k@mating.swarm_height_m\endcsname{2.87}
\expandafter\def\csname ms@k@mating.swarm_markers_per_km2\endcsname{1.0}
\expandafter\def\csname ms@k@mating.swarm_markers_source\endcsname{"engineering prior"}
\expandafter\def\csname ms@k@mating.swarm_start_min_after_sunset\endcsname{5.0}
\expandafter\def\csname ms@k@oviposition.detect_range_m\endcsname{2000.0}
\expandafter\def\csname ms@k@oviposition.min_attraction\endcsname{0.3}
\expandafter\def\csname ms@k@oviposition.scale_m\endcsname{50.0}
\expandafter\def\csname ms@k@oviposition.step_m\endcsname{50.0}
\expandafter\def\csname ms@k@oviposition.timeout_days\endcsname{7}
\expandafter\def\csname ms@k@pluvial.cold_start_z\endcsname{0.0}
\expandafter\def\csname ms@k@pluvial.delta_days\endcsname{1.0}
\expandafter\def\csname ms@k@pluvial.dry_memory_days\endcsname{5.0}
\expandafter\def\csname ms@k@pluvial.rain_scale_mm\endcsname{2.5}
\expandafter\def\csname ms@k@rng.aquatic_seed_larva_mortality\endcsname{123}
\expandafter\def\csname ms@k@rng.aquatic_seed_stage_mortality\endcsname{42}
\expandafter\def\csname ms@k@rng.aquatic_seed_washout\endcsname{789}
\expandafter\def\csname ms@k@species.arabiensis.activity_weights\endcsname{"0.02, 0.1, 0.8, 0.08"}
\expandafter\def\csname ms@k@species.arabiensis.aestivation_supported\endcsname{true}
\expandafter\def\csname ms@k@species.arabiensis.anemotaxis_deviation_deg\endcsname{30.0}
\expandafter\def\csname ms@k@species.arabiensis.anemotaxis_rule\endcsname{"upwind"}
\expandafter\def\csname ms@k@species.arabiensis.arrested_daily_survival\endcsname{0.986}
\expandafter\def\csname ms@k@species.arabiensis.blood_meal_refractory_days\endcsname{0.5}
\expandafter\def\csname ms@k@species.arabiensis.blood_meal_size_ul\endcsname{0.3}
\expandafter\def\csname ms@k@species.arabiensis.blood_meal_target_ul\endcsname{0.5}
\expandafter\def\csname ms@k@species.arabiensis.casting_angle_deg\endcsname{90.0}
\expandafter\def\csname ms@k@species.arabiensis.casting_probability\endcsname{0.1}
\expandafter\def\csname ms@k@species.arabiensis.co2_activation_radius_m\endcsname{100.0}
\expandafter\def\csname ms@k@species.arabiensis.daily_move_capacity_m\endcsname{590.0}
\expandafter\def\csname ms@k@species.arabiensis.directed_flight_budget_h\endcsname{0.4856}
\expandafter\def\csname ms@k@species.arabiensis.eggs_per_ul\endcsname{223.0}
\expandafter\def\csname ms@k@species.arabiensis.exploration_range_m\endcsname{100.0}
\expandafter\def\csname ms@k@species.arabiensis.exploration_step_correlation\endcsname{0.5}
\expandafter\def\csname ms@k@species.arabiensis.exploration_step_m\endcsname{50.0}
\expandafter\def\csname ms@k@species.arabiensis.first_cycle_target_ul\endcsname{0.8}
\expandafter\def\csname ms@k@species.arabiensis.flight_speed_ms\endcsname{0.3375}
\expandafter\def\csname ms@k@species.arabiensis.gonotrophic_arrest_prob_dry\endcsname{0.4454}
\expandafter\def\csname ms@k@species.arabiensis.male_daily_survival\endcsname{0.773}
\expandafter\def\csname ms@k@species.arabiensis.male_ratio\endcsname{0.5}
\expandafter\def\csname ms@k@species.arabiensis.meal_size_min_ul\endcsname{0.1}
\expandafter\def\csname ms@k@species.arabiensis.olfactory_radius_m\endcsname{50.0}
\expandafter\def\csname ms@k@species.arabiensis.plume_memory_s\endcsname{0.7}
\expandafter\def\csname ms@k@species.arabiensis.pref.cattle\endcsname{0.2}
\expandafter\def\csname ms@k@species.arabiensis.pref.goat\endcsname{0.15}
\expandafter\def\csname ms@k@species.arabiensis.pref.human\endcsname{0.5}
\expandafter\def\csname ms@k@species.arabiensis.pref.sheep\endcsname{0.1}
\expandafter\def\csname ms@k@species.arabiensis.pref.wildlife\endcsname{0.05}
\expandafter\def\csname ms@k@species.arabiensis.probing_success_rate\endcsname{0.7275}
\expandafter\def\csname ms@k@species.arabiensis.salinity_hi_tol_ppt\endcsname{15.0}
\expandafter\def\csname ms@k@species.arabiensis.salinity_opt_ppt\endcsname{0.0}
\expandafter\def\csname ms@k@species.arabiensis.salinity_slope\endcsname{0.15}
\expandafter\def\csname ms@k@species.arabiensis.terminal_radius_m\endcsname{0.15}
\expandafter\def\csname ms@k@species.arabiensis.visual_radius_m\endcsname{0.485}
\expandafter\def\csname ms@k@species.arabiensis.weights.crop_pluvial\endcsname{0.1}
\expandafter\def\csname ms@k@species.arabiensis.weights.open_pluvial\endcsname{0.55}
\expandafter\def\csname ms@k@species.arabiensis.weights.surface_persistent\endcsname{0.1}
\expandafter\def\csname ms@k@species.arabiensis.weights.urban_persistent\endcsname{0.15}
\expandafter\def\csname ms@k@species.arabiensis.weights.wetland_persistent\endcsname{0.1}
\expandafter\def\csname ms@k@species.arabiensis.weights.wetland_pluvial\endcsname{0.0}
\expandafter\def\csname ms@k@species.arabiensis.wind_block_threshold_ms\endcsname{0.3375}
\expandafter\def\csname ms@k@species.arabiensis.wind_reference\endcsname{"u100/v100"}
\expandafter\def\csname ms@k@species.coluzzii.activity_weights\endcsname{"0.0, 0.255, 0.48, 0.265"}
\expandafter\def\csname ms@k@species.coluzzii.aestivation_supported\endcsname{true}
\expandafter\def\csname ms@k@species.coluzzii.anemotaxis_deviation_deg\endcsname{30.0}
\expandafter\def\csname ms@k@species.coluzzii.anemotaxis_rule\endcsname{"upwind"}
\expandafter\def\csname ms@k@species.coluzzii.arrested_daily_survival\endcsname{0.986}
\expandafter\def\csname ms@k@species.coluzzii.blood_meal_refractory_days\endcsname{0.5}
\expandafter\def\csname ms@k@species.coluzzii.blood_meal_size_ul\endcsname{0.3}
\expandafter\def\csname ms@k@species.coluzzii.blood_meal_target_ul\endcsname{0.5}
\expandafter\def\csname ms@k@species.coluzzii.casting_angle_deg\endcsname{90.0}
\expandafter\def\csname ms@k@species.coluzzii.casting_probability\endcsname{0.1}
\expandafter\def\csname ms@k@species.coluzzii.co2_activation_radius_m\endcsname{100.0}
\expandafter\def\csname ms@k@species.coluzzii.daily_move_capacity_m\endcsname{590.0}
\expandafter\def\csname ms@k@species.coluzzii.directed_flight_budget_h\endcsname{0.4856}
\expandafter\def\csname ms@k@species.coluzzii.eggs_per_ul\endcsname{223.0}
\expandafter\def\csname ms@k@species.coluzzii.exploration_range_m\endcsname{100.0}
\expandafter\def\csname ms@k@species.coluzzii.exploration_step_correlation\endcsname{0.5}
\expandafter\def\csname ms@k@species.coluzzii.exploration_step_m\endcsname{50.0}
\expandafter\def\csname ms@k@species.coluzzii.first_cycle_target_ul\endcsname{0.8}
\expandafter\def\csname ms@k@species.coluzzii.flight_speed_ms\endcsname{0.3375}
\expandafter\def\csname ms@k@species.coluzzii.gonotrophic_arrest_prob_dry\endcsname{0.2704}
\expandafter\def\csname ms@k@species.coluzzii.male_daily_survival\endcsname{0.773}
\expandafter\def\csname ms@k@species.coluzzii.male_ratio\endcsname{0.5}
\expandafter\def\csname ms@k@species.coluzzii.meal_size_min_ul\endcsname{0.1}
\expandafter\def\csname ms@k@species.coluzzii.olfactory_radius_m\endcsname{50.0}
\expandafter\def\csname ms@k@species.coluzzii.plume_memory_s\endcsname{0.7}
\expandafter\def\csname ms@k@species.coluzzii.pref.cattle\endcsname{0.005}
\expandafter\def\csname ms@k@species.coluzzii.pref.goat\endcsname{0.003}
\expandafter\def\csname ms@k@species.coluzzii.pref.human\endcsname{0.99}
\expandafter\def\csname ms@k@species.coluzzii.pref.sheep\endcsname{0.001}
\expandafter\def\csname ms@k@species.coluzzii.pref.wildlife\endcsname{0.001}
\expandafter\def\csname ms@k@species.coluzzii.probing_success_rate\endcsname{0.7275}
\expandafter\def\csname ms@k@species.coluzzii.salinity_hi_tol_ppt\endcsname{30.0}
\expandafter\def\csname ms@k@species.coluzzii.salinity_opt_ppt\endcsname{0.0}
\expandafter\def\csname ms@k@species.coluzzii.salinity_slope\endcsname{0.08}
\expandafter\def\csname ms@k@species.coluzzii.terminal_radius_m\endcsname{0.15}
\expandafter\def\csname ms@k@species.coluzzii.visual_radius_m\endcsname{0.485}
\expandafter\def\csname ms@k@species.coluzzii.weights.crop_pluvial\endcsname{0.05}
\expandafter\def\csname ms@k@species.coluzzii.weights.open_pluvial\endcsname{0.25}
\expandafter\def\csname ms@k@species.coluzzii.weights.surface_persistent\endcsname{0.4}
\expandafter\def\csname ms@k@species.coluzzii.weights.urban_persistent\endcsname{0.2}
\expandafter\def\csname ms@k@species.coluzzii.weights.wetland_persistent\endcsname{0.1}
\expandafter\def\csname ms@k@species.coluzzii.weights.wetland_pluvial\endcsname{0.0}
\expandafter\def\csname ms@k@species.coluzzii.wind_block_threshold_ms\endcsname{0.3375}
\expandafter\def\csname ms@k@species.coluzzii.wind_reference\endcsname{"u100/v100"}
\expandafter\def\csname ms@k@species.default_species\endcsname{"coluzzii"}
\expandafter\def\csname ms@k@species.funestus.activity_weights\endcsname{"0.02, 0.1, 0.8, 0.08"}
\expandafter\def\csname ms@k@species.funestus.aestivation_supported\endcsname{false}
\expandafter\def\csname ms@k@species.funestus.anemotaxis_deviation_deg\endcsname{30.0}
\expandafter\def\csname ms@k@species.funestus.anemotaxis_rule\endcsname{"upwind"}
\expandafter\def\csname ms@k@species.funestus.arrested_daily_survival\endcsname{0.986}
\expandafter\def\csname ms@k@species.funestus.blood_meal_refractory_days\endcsname{0.5}
\expandafter\def\csname ms@k@species.funestus.blood_meal_size_ul\endcsname{0.3}
\expandafter\def\csname ms@k@species.funestus.blood_meal_target_ul\endcsname{0.5}
\expandafter\def\csname ms@k@species.funestus.casting_angle_deg\endcsname{90.0}
\expandafter\def\csname ms@k@species.funestus.casting_probability\endcsname{0.1}
\expandafter\def\csname ms@k@species.funestus.co2_activation_radius_m\endcsname{100.0}
\expandafter\def\csname ms@k@species.funestus.daily_move_capacity_m\endcsname{590.0}
\expandafter\def\csname ms@k@species.funestus.directed_flight_budget_h\endcsname{0.4856}
\expandafter\def\csname ms@k@species.funestus.eggs_per_ul\endcsname{223.0}
\expandafter\def\csname ms@k@species.funestus.exploration_range_m\endcsname{100.0}
\expandafter\def\csname ms@k@species.funestus.exploration_step_correlation\endcsname{0.5}
\expandafter\def\csname ms@k@species.funestus.exploration_step_m\endcsname{50.0}
\expandafter\def\csname ms@k@species.funestus.first_cycle_target_ul\endcsname{0.8}
\expandafter\def\csname ms@k@species.funestus.flight_speed_ms\endcsname{0.3375}
\expandafter\def\csname ms@k@species.funestus.gonotrophic_arrest_prob_dry\endcsname{0.0}
\expandafter\def\csname ms@k@species.funestus.male_daily_survival\endcsname{0.773}
\expandafter\def\csname ms@k@species.funestus.male_ratio\endcsname{0.5}
\expandafter\def\csname ms@k@species.funestus.meal_size_min_ul\endcsname{0.1}
\expandafter\def\csname ms@k@species.funestus.olfactory_radius_m\endcsname{50.0}
\expandafter\def\csname ms@k@species.funestus.plume_memory_s\endcsname{0.7}
\expandafter\def\csname ms@k@species.funestus.pref.cattle\endcsname{0.005}
\expandafter\def\csname ms@k@species.funestus.pref.goat\endcsname{0.003}
\expandafter\def\csname ms@k@species.funestus.pref.human\endcsname{0.99}
\expandafter\def\csname ms@k@species.funestus.pref.sheep\endcsname{0.001}
\expandafter\def\csname ms@k@species.funestus.pref.wildlife\endcsname{0.001}
\expandafter\def\csname ms@k@species.funestus.probing_success_rate\endcsname{0.7275}
\expandafter\def\csname ms@k@species.funestus.salinity_hi_tol_ppt\endcsname{15.0}
\expandafter\def\csname ms@k@species.funestus.salinity_opt_ppt\endcsname{0.0}
\expandafter\def\csname ms@k@species.funestus.salinity_slope\endcsname{0.15}
\expandafter\def\csname ms@k@species.funestus.terminal_radius_m\endcsname{0.15}
\expandafter\def\csname ms@k@species.funestus.visual_radius_m\endcsname{0.485}
\expandafter\def\csname ms@k@species.funestus.weights.crop_pluvial\endcsname{0.05}
\expandafter\def\csname ms@k@species.funestus.weights.open_pluvial\endcsname{0.1}
\expandafter\def\csname ms@k@species.funestus.weights.surface_persistent\endcsname{0.45}
\expandafter\def\csname ms@k@species.funestus.weights.urban_persistent\endcsname{0.1}
\expandafter\def\csname ms@k@species.funestus.weights.wetland_persistent\endcsname{0.3}
\expandafter\def\csname ms@k@species.funestus.weights.wetland_pluvial\endcsname{0.0}
\expandafter\def\csname ms@k@species.funestus.wind_block_threshold_ms\endcsname{0.3375}
\expandafter\def\csname ms@k@species.funestus.wind_reference\endcsname{"u100/v100"}
\expandafter\def\csname ms@k@species.gambiae_ss.activity_weights\endcsname{"0.02, 0.1, 0.8, 0.08"}
\expandafter\def\csname ms@k@species.gambiae_ss.aestivation_supported\endcsname{true}
\expandafter\def\csname ms@k@species.gambiae_ss.anemotaxis_deviation_deg\endcsname{30.0}
\expandafter\def\csname ms@k@species.gambiae_ss.anemotaxis_rule\endcsname{"upwind"}
\expandafter\def\csname ms@k@species.gambiae_ss.arrested_daily_survival\endcsname{0.986}
\expandafter\def\csname ms@k@species.gambiae_ss.blood_meal_refractory_days\endcsname{0.5}
\expandafter\def\csname ms@k@species.gambiae_ss.blood_meal_size_ul\endcsname{0.3}
\expandafter\def\csname ms@k@species.gambiae_ss.blood_meal_target_ul\endcsname{0.5}
\expandafter\def\csname ms@k@species.gambiae_ss.casting_angle_deg\endcsname{90.0}
\expandafter\def\csname ms@k@species.gambiae_ss.casting_probability\endcsname{0.1}
\expandafter\def\csname ms@k@species.gambiae_ss.co2_activation_radius_m\endcsname{100.0}
\expandafter\def\csname ms@k@species.gambiae_ss.daily_move_capacity_m\endcsname{590.0}
\expandafter\def\csname ms@k@species.gambiae_ss.directed_flight_budget_h\endcsname{0.4856}
\expandafter\def\csname ms@k@species.gambiae_ss.eggs_per_ul\endcsname{223.0}
\expandafter\def\csname ms@k@species.gambiae_ss.exploration_range_m\endcsname{100.0}
\expandafter\def\csname ms@k@species.gambiae_ss.exploration_step_correlation\endcsname{0.5}
\expandafter\def\csname ms@k@species.gambiae_ss.exploration_step_m\endcsname{50.0}
\expandafter\def\csname ms@k@species.gambiae_ss.first_cycle_target_ul\endcsname{0.8}
\expandafter\def\csname ms@k@species.gambiae_ss.flight_speed_ms\endcsname{0.3375}
\expandafter\def\csname ms@k@species.gambiae_ss.gonotrophic_arrest_prob_dry\endcsname{0.6838}
\expandafter\def\csname ms@k@species.gambiae_ss.male_daily_survival\endcsname{0.773}
\expandafter\def\csname ms@k@species.gambiae_ss.male_ratio\endcsname{0.5}
\expandafter\def\csname ms@k@species.gambiae_ss.meal_size_min_ul\endcsname{0.1}
\expandafter\def\csname ms@k@species.gambiae_ss.olfactory_radius_m\endcsname{50.0}
\expandafter\def\csname ms@k@species.gambiae_ss.plume_memory_s\endcsname{0.7}
\expandafter\def\csname ms@k@species.gambiae_ss.pref.cattle\endcsname{0.005}
\expandafter\def\csname ms@k@species.gambiae_ss.pref.goat\endcsname{0.003}
\expandafter\def\csname ms@k@species.gambiae_ss.pref.human\endcsname{0.99}
\expandafter\def\csname ms@k@species.gambiae_ss.pref.sheep\endcsname{0.001}
\expandafter\def\csname ms@k@species.gambiae_ss.pref.wildlife\endcsname{0.001}
\expandafter\def\csname ms@k@species.gambiae_ss.probing_success_rate\endcsname{0.7275}
\expandafter\def\csname ms@k@species.gambiae_ss.salinity_hi_tol_ppt\endcsname{20.0}
\expandafter\def\csname ms@k@species.gambiae_ss.salinity_opt_ppt\endcsname{0.0}
\expandafter\def\csname ms@k@species.gambiae_ss.salinity_slope\endcsname{0.12}
\expandafter\def\csname ms@k@species.gambiae_ss.terminal_radius_m\endcsname{0.15}
\expandafter\def\csname ms@k@species.gambiae_ss.visual_radius_m\endcsname{0.485}
\expandafter\def\csname ms@k@species.gambiae_ss.weights.crop_pluvial\endcsname{0.1}
\expandafter\def\csname ms@k@species.gambiae_ss.weights.open_pluvial\endcsname{0.65}
\expandafter\def\csname ms@k@species.gambiae_ss.weights.surface_persistent\endcsname{0.05}
\expandafter\def\csname ms@k@species.gambiae_ss.weights.urban_persistent\endcsname{0.15}
\expandafter\def\csname ms@k@species.gambiae_ss.weights.wetland_persistent\endcsname{0.05}
\expandafter\def\csname ms@k@species.gambiae_ss.weights.wetland_pluvial\endcsname{0.0}
\expandafter\def\csname ms@k@species.gambiae_ss.wind_block_threshold_ms\endcsname{0.3375}
\expandafter\def\csname ms@k@species.gambiae_ss.wind_reference\endcsname{"u100/v100"}
\expandafter\def\csname ms@k@species.k_ref\endcsname{50.0}
\expandafter\def\csname ms@k@species.melas.activity_weights\endcsname{"0.02, 0.1, 0.8, 0.08"}
\expandafter\def\csname ms@k@species.melas.aestivation_supported\endcsname{false}
\expandafter\def\csname ms@k@species.melas.anemotaxis_deviation_deg\endcsname{30.0}
\expandafter\def\csname ms@k@species.melas.anemotaxis_rule\endcsname{"upwind"}
\expandafter\def\csname ms@k@species.melas.arrested_daily_survival\endcsname{0.986}
\expandafter\def\csname ms@k@species.melas.blood_meal_refractory_days\endcsname{0.5}
\expandafter\def\csname ms@k@species.melas.blood_meal_size_ul\endcsname{0.3}
\expandafter\def\csname ms@k@species.melas.blood_meal_target_ul\endcsname{0.5}
\expandafter\def\csname ms@k@species.melas.casting_angle_deg\endcsname{90.0}
\expandafter\def\csname ms@k@species.melas.casting_probability\endcsname{0.1}
\expandafter\def\csname ms@k@species.melas.co2_activation_radius_m\endcsname{100.0}
\expandafter\def\csname ms@k@species.melas.daily_move_capacity_m\endcsname{590.0}
\expandafter\def\csname ms@k@species.melas.directed_flight_budget_h\endcsname{0.4856}
\expandafter\def\csname ms@k@species.melas.eggs_per_ul\endcsname{223.0}
\expandafter\def\csname ms@k@species.melas.exploration_range_m\endcsname{100.0}
\expandafter\def\csname ms@k@species.melas.exploration_step_correlation\endcsname{0.5}
\expandafter\def\csname ms@k@species.melas.exploration_step_m\endcsname{50.0}
\expandafter\def\csname ms@k@species.melas.first_cycle_target_ul\endcsname{0.8}
\expandafter\def\csname ms@k@species.melas.flight_speed_ms\endcsname{0.3375}
\expandafter\def\csname ms@k@species.melas.gonotrophic_arrest_prob_dry\endcsname{0.0}
\expandafter\def\csname ms@k@species.melas.male_daily_survival\endcsname{0.773}
\expandafter\def\csname ms@k@species.melas.male_ratio\endcsname{0.5}
\expandafter\def\csname ms@k@species.melas.meal_size_min_ul\endcsname{0.1}
\expandafter\def\csname ms@k@species.melas.olfactory_radius_m\endcsname{50.0}
\expandafter\def\csname ms@k@species.melas.plume_memory_s\endcsname{0.7}
\expandafter\def\csname ms@k@species.melas.pref.cattle\endcsname{0.005}
\expandafter\def\csname ms@k@species.melas.pref.goat\endcsname{0.003}
\expandafter\def\csname ms@k@species.melas.pref.human\endcsname{0.99}
\expandafter\def\csname ms@k@species.melas.pref.sheep\endcsname{0.001}
\expandafter\def\csname ms@k@species.melas.pref.wildlife\endcsname{0.001}
\expandafter\def\csname ms@k@species.melas.probing_success_rate\endcsname{0.7275}
\expandafter\def\csname ms@k@species.melas.salinity_hi_tol_ppt\endcsname{38.0}
\expandafter\def\csname ms@k@species.melas.salinity_opt_ppt\endcsname{0.0}
\expandafter\def\csname ms@k@species.melas.salinity_slope\endcsname{0.02}
\expandafter\def\csname ms@k@species.melas.terminal_radius_m\endcsname{0.15}
\expandafter\def\csname ms@k@species.melas.visual_radius_m\endcsname{0.485}
\expandafter\def\csname ms@k@species.melas.weights.crop_pluvial\endcsname{0.05}
\expandafter\def\csname ms@k@species.melas.weights.open_pluvial\endcsname{0.4}
\expandafter\def\csname ms@k@species.melas.weights.surface_persistent\endcsname{0.25}
\expandafter\def\csname ms@k@species.melas.weights.urban_persistent\endcsname{0.15}
\expandafter\def\csname ms@k@species.melas.weights.wetland_persistent\endcsname{0.15}
\expandafter\def\csname ms@k@species.melas.weights.wetland_pluvial\endcsname{0.0}
\expandafter\def\csname ms@k@species.melas.wind_block_threshold_ms\endcsname{0.3375}
\expandafter\def\csname ms@k@species.melas.wind_reference\endcsname{"u100/v100"}
\expandafter\def\csname ms@k@species.neutral.weights.crop_pluvial\endcsname{0.05}
\expandafter\def\csname ms@k@species.neutral.weights.open_pluvial\endcsname{0.4}
\expandafter\def\csname ms@k@species.neutral.weights.surface_persistent\endcsname{0.25}
\expandafter\def\csname ms@k@species.neutral.weights.urban_persistent\endcsname{0.15}
\expandafter\def\csname ms@k@species.neutral.weights.wetland_persistent\endcsname{0.15}
\expandafter\def\csname ms@k@species.neutral.weights.wetland_pluvial\endcsname{0.0}
\expandafter\def\csname ms@k@species.stephensi.activity_weights\endcsname{"0.02, 0.1, 0.8, 0.08"}
\expandafter\def\csname ms@k@species.stephensi.aestivation_supported\endcsname{false}
\expandafter\def\csname ms@k@species.stephensi.anemotaxis_deviation_deg\endcsname{30.0}
\expandafter\def\csname ms@k@species.stephensi.anemotaxis_rule\endcsname{"upwind"}
\expandafter\def\csname ms@k@species.stephensi.arrested_daily_survival\endcsname{0.986}
\expandafter\def\csname ms@k@species.stephensi.blood_meal_refractory_days\endcsname{0.5}
\expandafter\def\csname ms@k@species.stephensi.blood_meal_size_ul\endcsname{0.3}
\expandafter\def\csname ms@k@species.stephensi.blood_meal_target_ul\endcsname{0.5}
\expandafter\def\csname ms@k@species.stephensi.casting_angle_deg\endcsname{90.0}
\expandafter\def\csname ms@k@species.stephensi.casting_probability\endcsname{0.1}
\expandafter\def\csname ms@k@species.stephensi.co2_activation_radius_m\endcsname{100.0}
\expandafter\def\csname ms@k@species.stephensi.daily_move_capacity_m\endcsname{590.0}
\expandafter\def\csname ms@k@species.stephensi.directed_flight_budget_h\endcsname{0.4856}
\expandafter\def\csname ms@k@species.stephensi.eggs_per_ul\endcsname{223.0}
\expandafter\def\csname ms@k@species.stephensi.exploration_range_m\endcsname{100.0}
\expandafter\def\csname ms@k@species.stephensi.exploration_step_correlation\endcsname{0.5}
\expandafter\def\csname ms@k@species.stephensi.exploration_step_m\endcsname{50.0}
\expandafter\def\csname ms@k@species.stephensi.first_cycle_target_ul\endcsname{0.8}
\expandafter\def\csname ms@k@species.stephensi.flight_speed_ms\endcsname{0.3375}
\expandafter\def\csname ms@k@species.stephensi.gonotrophic_arrest_prob_dry\endcsname{0.0}
\expandafter\def\csname ms@k@species.stephensi.male_daily_survival\endcsname{0.773}
\expandafter\def\csname ms@k@species.stephensi.male_ratio\endcsname{0.5}
\expandafter\def\csname ms@k@species.stephensi.meal_size_min_ul\endcsname{0.1}
\expandafter\def\csname ms@k@species.stephensi.olfactory_radius_m\endcsname{50.0}
\expandafter\def\csname ms@k@species.stephensi.plume_memory_s\endcsname{0.7}
\expandafter\def\csname ms@k@species.stephensi.pref.cattle\endcsname{0.15}
\expandafter\def\csname ms@k@species.stephensi.pref.goat\endcsname{0.05}
\expandafter\def\csname ms@k@species.stephensi.pref.human\endcsname{0.7}
\expandafter\def\csname ms@k@species.stephensi.pref.sheep\endcsname{0.05}
\expandafter\def\csname ms@k@species.stephensi.pref.wildlife\endcsname{0.05}
\expandafter\def\csname ms@k@species.stephensi.probing_success_rate\endcsname{0.7275}
\expandafter\def\csname ms@k@species.stephensi.salinity_hi_tol_ppt\endcsname{20.0}
\expandafter\def\csname ms@k@species.stephensi.salinity_opt_ppt\endcsname{0.0}
\expandafter\def\csname ms@k@species.stephensi.salinity_slope\endcsname{0.15}
\expandafter\def\csname ms@k@species.stephensi.terminal_radius_m\endcsname{0.15}
\expandafter\def\csname ms@k@species.stephensi.visual_radius_m\endcsname{0.485}
\expandafter\def\csname ms@k@species.stephensi.weights.crop_pluvial\endcsname{0.05}
\expandafter\def\csname ms@k@species.stephensi.weights.open_pluvial\endcsname{0.4}
\expandafter\def\csname ms@k@species.stephensi.weights.surface_persistent\endcsname{0.25}
\expandafter\def\csname ms@k@species.stephensi.weights.urban_persistent\endcsname{0.15}
\expandafter\def\csname ms@k@species.stephensi.weights.wetland_persistent\endcsname{0.15}
\expandafter\def\csname ms@k@species.stephensi.weights.wetland_pluvial\endcsname{0.0}
\expandafter\def\csname ms@k@species.stephensi.wind_block_threshold_ms\endcsname{0.3375}
\expandafter\def\csname ms@k@species.stephensi.wind_reference\endcsname{"u100/v100"}
\expandafter\def\csname ms@k@thermal.efd_max_eggs_per_day\endcsname{0.5}
\expandafter\def\csname ms@k@thermal.efd_tmax_c\endcsname{34.0}
\expandafter\def\csname ms@k@thermal.efd_tmin_c\endcsname{16.0}
\expandafter\def\csname ms@k@thermal.efd_topt_c\endcsname{25.0}
\expandafter\def\csname ms@k@thermal.egg_development_scale\endcsname{0.00062}
\expandafter\def\csname ms@k@thermal.egg_development_tmax_c\endcsname{40.0}
\expandafter\def\csname ms@k@thermal.egg_development_tmin_c\endcsname{10.0}
\expandafter\def\csname ms@k@thermal.larva_development_scale\endcsname{0.001}
\expandafter\def\csname ms@k@thermal.larva_development_tmax_c\endcsname{40.0}
\expandafter\def\csname ms@k@thermal.larva_development_tmin_c\endcsname{10.0}
\expandafter\def\csname ms@k@thermal.pupa_development_scale\endcsname{0.00066}
\expandafter\def\csname ms@k@thermal.pupa_development_tmax_c\endcsname{40.0}
\expandafter\def\csname ms@k@thermal.pupa_development_tmin_c\endcsname{12.0}
\expandafter\def\csname ms@k@transmission.beta_hv\endcsname{0.4}
\expandafter\def\csname ms@k@transmission.beta_vh\endcsname{0.5}
\expandafter\def\csname ms@k@transmission.human_immunity_days\endcsname{180}
\expandafter\def\csname ms@k@transmission.human_incubation_days\endcsname{12}
\expandafter\def\csname ms@k@transmission.human_infectious_days\endcsname{20}
\expandafter\def\csname ms@k@transmission.initial_human_prevalence\endcsname{0.05}
\expandafter\def\csname ms@k@transmission.initial_vector_infected_frac\endcsname{0.0}
\expandafter\def\csname ms@k@urban.beta_urban\endcsname{0.75}
\expandafter\def\csname ms@k@urban.rain_pools_built_damage\endcsname{0.3}
\expandafter\def\csname ms@k@version\endcsname{3}
\expandafter\def\csname ms@k@wildlife.class_suitability.bare\endcsname{0.2}
\expandafter\def\csname ms@k@wildlife.class_suitability.built\endcsname{0.05}
\expandafter\def\csname ms@k@wildlife.class_suitability.crops\endcsname{0.3}
\expandafter\def\csname ms@k@wildlife.class_suitability.flooded_vegetation\endcsname{0.45}
\expandafter\def\csname ms@k@wildlife.class_suitability.grass\endcsname{0.6}
\expandafter\def\csname ms@k@wildlife.class_suitability.shrub_and_scrub\endcsname{0.7}
\expandafter\def\csname ms@k@wildlife.class_suitability.snow_and_ice\endcsname{0.1}
\expandafter\def\csname ms@k@wildlife.class_suitability.trees\endcsname{0.9}
\expandafter\def\csname ms@k@wildlife.class_suitability.water\endcsname{0.0}
\expandafter\def\csname ms@k@wildlife.weights.habitat\endcsname{0.5}
\expandafter\def\csname ms@k@wildlife.weights.remoteness\endcsname{0.2}
\expandafter\def\csname ms@k@wildlife.weights.water\endcsname{0.3}

% Generated by tools/gen_concepts_values.py -- do not edit.
% Values are read from the CODE that implements each stage, by AST.
% A key that no longer exists makes \codeconst fail the LaTeX build
% with an error, in any document that references it.
% Each line carries the file:line the value was read from.

\makeatletter
\newcommand{\codeconst}[1]{%
  \expandafter\ifx\csname cc@k@#1\endcsname\relax
    \errmessage{Concept document: unknown code constant `#1'}%
  \else
    \csname cc@k@#1\endcsname%
  \fi}
\makeatother

% download.wildlife_habitat_fallback = 0.5  (mal-commonlib/src/mal_commonlib/data/loaders/wildlife.py:205)
\expandafter\def\csname cc@k@download.wildlife_habitat_fallback\endcsname{0.5}
% download.wildlife_remoteness_neutral = 1.0  (mal-commonlib/src/mal_commonlib/data/loaders/wildlife.py:161)
\expandafter\def\csname cc@k@download.wildlife_remoteness_neutral\endcsname{1.0}
% ingest.bbox_east = 1.5  (mal-commonlib/src/mal_commonlib/aois.yaml:18)
\expandafter\def\csname cc@k@ingest.bbox_east\endcsname{1.5}
% ingest.bbox_north = 11.5  (mal-commonlib/src/mal_commonlib/aois.yaml:18)
\expandafter\def\csname cc@k@ingest.bbox_north\endcsname{11.5}
% ingest.bbox_south = 4.5  (mal-commonlib/src/mal_commonlib/aois.yaml:18)
\expandafter\def\csname cc@k@ingest.bbox_south\endcsname{4.5}
% ingest.bbox_west = -3.5  (mal-commonlib/src/mal_commonlib/aois.yaml:18)
\expandafter\def\csname cc@k@ingest.bbox_west\endcsname{-3.5}
% ingest.cell_size_m = 1000  (mal-commonlib/src/mal_commonlib/aois.yaml:20)
\expandafter\def\csname cc@k@ingest.cell_size_m\endcsname{1000}
% ingest.env_daily_channels = 2  (mal-core/src/mal_core/ingest/daily_nc.py:414)
\expandafter\def\csname cc@k@ingest.env_daily_channels\endcsname{2}
% ingest.env_static_planes = 5  (mal-core/src/mal_core/ingest/daily_nc.py:414)
\expandafter\def\csname cc@k@ingest.env_static_planes\endcsname{5}
% ingest.fallback_building_rural = 0.1  (mal-core/src/mal_core/ingest/hosts.py:142)
\expandafter\def\csname cc@k@ingest.fallback_building_rural\endcsname{0.1}
% ingest.fallback_building_urban = 0.6  (mal-core/src/mal_core/ingest/hosts.py:142)
\expandafter\def\csname cc@k@ingest.fallback_building_urban\endcsname{0.6}
% ingest.fallback_wildlife_proxy = 0.3  (mal-core/src/mal_core/ingest/hosts.py:165)
\expandafter\def\csname cc@k@ingest.fallback_wildlife_proxy\endcsname{0.3}
% ingest.jrc_occurrence_threshold_pct = 80  (mal-commonlib/src/mal_commonlib/data/loaders/jrc_gsw.py:72)
\expandafter\def\csname cc@k@ingest.jrc_occurrence_threshold_pct\endcsname{80}
% ingest.k_ref_diag_per_km2 = 1329900.0  (mal-core/src/mal_core/habitat/capacity.py:30)
\expandafter\def\csname cc@k@ingest.k_ref_diag_per_km2\endcsname{1329900.0}
% ingest.mobility_beta_day = 0.05  (mal-core/src/mal_core/ingest/mobility.py:43)
\expandafter\def\csname cc@k@ingest.mobility_beta_day\endcsname{0.05}
% ingest.mobility_beta_livestock = 0.1  (mal-core/src/mal_core/ingest/mobility.py:45)
\expandafter\def\csname cc@k@ingest.mobility_beta_livestock\endcsname{0.1}
% ingest.mobility_beta_night = 0.5  (mal-core/src/mal_core/ingest/mobility.py:44)
\expandafter\def\csname cc@k@ingest.mobility_beta_night\endcsname{0.5}
% ingest.mobility_cell_size_km = 1.0  (mal-core/src/mal_core/ingest/mobility.py:42)
\expandafter\def\csname cc@k@ingest.mobility_cell_size_km\endcsname{1.0}
% ingest.mobility_max_distance_km = 50.0  (mal-core/src/mal_core/ingest/mobility.py:46)
\expandafter\def\csname cc@k@ingest.mobility_max_distance_km\endcsname{50.0}
% ingest.mobility_sparsity_threshold = 0.0005  (mal-core/src/mal_core/ingest/mobility.py:47)
\expandafter\def\csname cc@k@ingest.mobility_sparsity_threshold\endcsname{0.0005}
% ingest.urban_class_rural = 50  (mal-core/src/mal_core/ingest/hosts.py:126)
\expandafter\def\csname cc@k@ingest.urban_class_rural\endcsname{50}
% ingest.urban_class_urban = 30  (mal-core/src/mal_core/ingest/hosts.py:125)
\expandafter\def\csname cc@k@ingest.urban_class_urban\endcsname{30}
% ingest.urban_class_water = 20  (mal-core/src/mal_core/ingest/hosts.py:127)
\expandafter\def\csname cc@k@ingest.urban_class_water\endcsname{20}
% prediction.block_size = 128  (mal-core/src/mal_core/prediction/predictor.py:126)
\expandafter\def\csname cc@k@prediction.block_size\endcsname{128}
% prediction.contract_version_manifest = "1.1"  (mal-core/src/mal_core/prediction/registry.py:31)
\expandafter\def\csname cc@k@prediction.contract_version_manifest\endcsname{"1.1"}
% prediction.contract_version_sidecar = "1.1"  (mal-core/src/mal_core/prediction/predictor.py:146)
\expandafter\def\csname cc@k@prediction.contract_version_sidecar\endcsname{"1.1"}
% prediction.in_channels = 10  (mal-core/src/mal_core/prediction/registry.py:32)
\expandafter\def\csname cc@k@prediction.in_channels\endcsname{10}
% prediction.manifest_version = "1.0"  (mal-core/src/mal_core/prediction/registry.py:30)
\expandafter\def\csname cc@k@prediction.manifest_version\endcsname{"1.0"}
% prediction.out_channels = 6  (mal-core/src/mal_core/prediction/registry.py:33)
\expandafter\def\csname cc@k@prediction.out_channels\endcsname{6}
% prediction.scale_continental_m = 5000  (mal-core/src/mal_core/prediction/aggregator.py:29)
\expandafter\def\csname cc@k@prediction.scale_continental_m\endcsname{5000}
% prediction.scale_national_m = 2000  (mal-core/src/mal_core/prediction/aggregator.py:28)
\expandafter\def\csname cc@k@prediction.scale_national_m\endcsname{2000}
% prediction.scale_regional_m = 1000  (mal-core/src/mal_core/prediction/aggregator.py:27)
\expandafter\def\csname cc@k@prediction.scale_regional_m\endcsname{1000}
% scoring.d11_larva_band_max = 0.9  (mal-core/src/mal_core/scoring/scorers/d11_aquatic_structure.py:35)
\expandafter\def\csname cc@k@scoring.d11_larva_band_max\endcsname{0.9}
% scoring.d11_larva_band_min = 0.65  (mal-core/src/mal_core/scoring/scorers/d11_aquatic_structure.py:35)
\expandafter\def\csname cc@k@scoring.d11_larva_band_min\endcsname{0.65}
% scoring.d11_pupa_band_max = 0.15  (mal-core/src/mal_core/scoring/scorers/d11_aquatic_structure.py:36)
\expandafter\def\csname cc@k@scoring.d11_pupa_band_max\endcsname{0.15}
% scoring.d11_pupa_band_min = 0.02  (mal-core/src/mal_core/scoring/scorers/d11_aquatic_structure.py:36)
\expandafter\def\csname cc@k@scoring.d11_pupa_band_min\endcsname{0.02}
% scoring.d15_expansion_fraction = 0.2  (mal-core/src/mal_core/scoring/scorers/d15_persistence.py:125)
\expandafter\def\csname cc@k@scoring.d15_expansion_fraction\endcsname{0.2}
% scoring.d15_min_horizon_days = 30  (mal-core/src/mal_core/scoring/scorers/d15_persistence.py:97)
\expandafter\def\csname cc@k@scoring.d15_min_horizon_days\endcsname{30}
% scoring.d15_survival_fraction = 0.3  (mal-core/src/mal_core/scoring/scorers/d15_persistence.py:105)
\expandafter\def\csname cc@k@scoring.d15_survival_fraction\endcsname{0.3}
% scoring.d16_coverage_target = 0.5  (mal-core/src/mal_core/scoring/scorers/d16_detection_coverage.py:44)
\expandafter\def\csname cc@k@scoring.d16_coverage_target\endcsname{0.5}
% scoring.d16_radius_km_default = 5.0  (mal-core/src/mal_core/scoring/scorers/d16_detection_coverage.py:43)
\expandafter\def\csname cc@k@scoring.d16_radius_km_default\endcsname{5.0}
% scoring.d16_sensitivity_radius_far_km = 10.0  (mal-core/src/mal_core/scoring/scorers/d16_detection_coverage.py:45)
\expandafter\def\csname cc@k@scoring.d16_sensitivity_radius_far_km\endcsname{10.0}
% scoring.d16_sensitivity_radius_mid_km = 5.0  (mal-core/src/mal_core/scoring/scorers/d16_detection_coverage.py:45)
\expandafter\def\csname cc@k@scoring.d16_sensitivity_radius_mid_km\endcsname{5.0}
% scoring.d16_sensitivity_radius_near_km = 2.0  (mal-core/src/mal_core/scoring/scorers/d16_detection_coverage.py:45)
\expandafter\def\csname cc@k@scoring.d16_sensitivity_radius_near_km\endcsname{2.0}
% scoring.d19_activation_floor = 0.05  (mal-core/src/mal_core/scoring/scorers/d19_pool_persistence.py:81)
\expandafter\def\csname cc@k@scoring.d19_activation_floor\endcsname{0.05}
% scoring.d19_check_drying_weight = 0.2  (mal-core/src/mal_core/scoring/scorers/d19_pool_persistence.py:130)
\expandafter\def\csname cc@k@scoring.d19_check_drying_weight\endcsname{0.2}
% scoring.d19_check_floor_weight = 0.3  (mal-core/src/mal_core/scoring/scorers/d19_pool_persistence.py:134)
\expandafter\def\csname cc@k@scoring.d19_check_floor_weight\endcsname{0.3}
% scoring.d19_check_monotone_weight = 0.2  (mal-core/src/mal_core/scoring/scorers/d19_pool_persistence.py:132)
\expandafter\def\csname cc@k@scoring.d19_check_monotone_weight\endcsname{0.2}
% scoring.d19_check_persistence_weight = 0.3  (mal-core/src/mal_core/scoring/scorers/d19_pool_persistence.py:128)
\expandafter\def\csname cc@k@scoring.d19_check_persistence_weight\endcsname{0.3}
% scoring.d19_dry_days = 10  (mal-core/src/mal_core/scoring/scorers/d19_pool_persistence.py:84)
\expandafter\def\csname cc@k@scoring.d19_dry_days\endcsname{10}
% scoring.d19_pulse_days = 3  (mal-core/src/mal_core/scoring/scorers/d19_pool_persistence.py:78)
\expandafter\def\csname cc@k@scoring.d19_pulse_days\endcsname{3}
% scoring.d19_rain_pulse_mm = 30.0  (mal-core/src/mal_core/scoring/scorers/d19_pool_persistence.py:79)
\expandafter\def\csname cc@k@scoring.d19_rain_pulse_mm\endcsname{30.0}
% scoring.d1_band_max = 0.7  (mal-core/src/mal_core/scoring/scorers/d1_habitat_occupancy.py:33)
\expandafter\def\csname cc@k@scoring.d1_band_max\endcsname{0.7}
% scoring.d1_band_min = 0.2  (mal-core/src/mal_core/scoring/scorers/d1_habitat_occupancy.py:33)
\expandafter\def\csname cc@k@scoring.d1_band_min\endcsname{0.2}
% scoring.d1_center = 0.45  (mal-core/src/mal_core/scoring/scorers/d1_habitat_occupancy.py:34)
\expandafter\def\csname cc@k@scoring.d1_center\endcsname{0.45}
% scoring.d1_sigma = 0.15  (mal-core/src/mal_core/scoring/scorers/d1_habitat_occupancy.py:35)
\expandafter\def\csname cc@k@scoring.d1_sigma\endcsname{0.15}
% scoring.d20_cap_probe_mm = 80.0  (mal-core/src/mal_core/scoring/scorers/d20_washout_response.py:56)
\expandafter\def\csname cc@k@scoring.d20_cap_probe_mm\endcsname{80.0}
% scoring.d20_check_cap_weight = 0.2  (mal-core/src/mal_core/scoring/scorers/d20_washout_response.py:71)
\expandafter\def\csname cc@k@scoring.d20_check_cap_weight\endcsname{0.2}
% scoring.d20_check_no_washout_weight = 0.3  (mal-core/src/mal_core/scoring/scorers/d20_washout_response.py:69)
\expandafter\def\csname cc@k@scoring.d20_check_no_washout_weight\endcsname{0.3}
% scoring.d20_check_threshold_weight = 0.2  (mal-core/src/mal_core/scoring/scorers/d20_washout_response.py:73)
\expandafter\def\csname cc@k@scoring.d20_check_threshold_weight\endcsname{0.2}
% scoring.d20_check_washout_weight = 0.3  (mal-core/src/mal_core/scoring/scorers/d20_washout_response.py:67)
\expandafter\def\csname cc@k@scoring.d20_check_washout_weight\endcsname{0.3}
% scoring.d20_min_washout_fraction = 0.15  (mal-core/src/mal_core/scoring/scorers/d20_washout_response.py:48)
\expandafter\def\csname cc@k@scoring.d20_min_washout_fraction\endcsname{0.15}
% scoring.d20_no_washout_probe_mm = 30.0  (mal-core/src/mal_core/scoring/scorers/d20_washout_response.py:51)
\expandafter\def\csname cc@k@scoring.d20_no_washout_probe_mm\endcsname{30.0}
% scoring.d20_threshold_max_fraction = 0.01  (mal-core/src/mal_core/scoring/scorers/d20_washout_response.py:63)
\expandafter\def\csname cc@k@scoring.d20_threshold_max_fraction\endcsname{0.01}
% scoring.d20_threshold_probe_mm = 40.0  (mal-core/src/mal_core/scoring/scorers/d20_washout_response.py:61)
\expandafter\def\csname cc@k@scoring.d20_threshold_probe_mm\endcsname{40.0}
% scoring.d20_washout_probe_mm = 50.0  (mal-core/src/mal_core/scoring/scorers/d20_washout_response.py:46)
\expandafter\def\csname cc@k@scoring.d20_washout_probe_mm\endcsname{50.0}
% scoring.d21_horizon_day = 30  (mal-core/src/mal_core/scoring/scorers/d21_dispersal_kernel.py:112)
\expandafter\def\csname cc@k@scoring.d21_horizon_day\endcsname{30}
% scoring.d21_p50_band_max = 0.7  (mal-core/src/mal_core/scoring/scorers/d21_dispersal_kernel.py:39)
\expandafter\def\csname cc@k@scoring.d21_p50_band_max\endcsname{0.7}
% scoring.d21_p50_band_min = 0.15  (mal-core/src/mal_core/scoring/scorers/d21_dispersal_kernel.py:39)
\expandafter\def\csname cc@k@scoring.d21_p50_band_min\endcsname{0.15}
% scoring.d21_p50_center = 0.4  (mal-core/src/mal_core/scoring/scorers/d21_dispersal_kernel.py:41)
\expandafter\def\csname cc@k@scoring.d21_p50_center\endcsname{0.4}
% scoring.d21_p50_sigma = 0.2  (mal-core/src/mal_core/scoring/scorers/d21_dispersal_kernel.py:41)
\expandafter\def\csname cc@k@scoring.d21_p50_sigma\endcsname{0.2}
% scoring.d21_p90_band_max = 3.0  (mal-core/src/mal_core/scoring/scorers/d21_dispersal_kernel.py:40)
\expandafter\def\csname cc@k@scoring.d21_p90_band_max\endcsname{3.0}
% scoring.d21_p90_band_min = 0.7  (mal-core/src/mal_core/scoring/scorers/d21_dispersal_kernel.py:40)
\expandafter\def\csname cc@k@scoring.d21_p90_band_min\endcsname{0.7}
% scoring.d21_p90_center = 1.6  (mal-core/src/mal_core/scoring/scorers/d21_dispersal_kernel.py:42)
\expandafter\def\csname cc@k@scoring.d21_p90_center\endcsname{1.6}
% scoring.d21_p90_sigma = 0.8  (mal-core/src/mal_core/scoring/scorers/d21_dispersal_kernel.py:42)
\expandafter\def\csname cc@k@scoring.d21_p90_sigma\endcsname{0.8}
% scoring.d25_r_target = 0.6  (mal-core/src/mal_core/scoring/scorers/d25_cases_ghana.py:45)
\expandafter\def\csname cc@k@scoring.d25_r_target\endcsname{0.6}
% scoring.d2_lifetime_centre = 11.0  (mal-core/src/mal_core/scoring/scorers/d2_survival.py:70)
\expandafter\def\csname cc@k@scoring.d2_lifetime_centre\endcsname{11.0}
% scoring.d2_lifetime_width = 4.0  (mal-core/src/mal_core/scoring/scorers/d2_survival.py:70)
\expandafter\def\csname cc@k@scoring.d2_lifetime_width\endcsname{4.0}
% scoring.d2_window_max = 15  (mal-core/src/mal_core/scoring/scorers/d2_survival.py:47)
\expandafter\def\csname cc@k@scoring.d2_window_max\endcsname{15}
% scoring.d2_window_min = 8  (mal-core/src/mal_core/scoring/scorers/d2_survival.py:47)
\expandafter\def\csname cc@k@scoring.d2_window_min\endcsname{8}
% scoring.d3_band_max = 0.5  (mal-core/src/mal_core/scoring/scorers/d3_eip.py:69)
\expandafter\def\csname cc@k@scoring.d3_band_max\endcsname{0.5}
% scoring.d3_band_min = 0.2  (mal-core/src/mal_core/scoring/scorers/d3_eip.py:69)
\expandafter\def\csname cc@k@scoring.d3_band_min\endcsname{0.2}
% scoring.d3_center = 0.35  (mal-core/src/mal_core/scoring/scorers/d3_eip.py:72)
\expandafter\def\csname cc@k@scoring.d3_center\endcsname{0.35}
% scoring.d3_day = 30  (mal-core/src/mal_core/scoring/scorers/d3_eip.py:67)
\expandafter\def\csname cc@k@scoring.d3_day\endcsname{30}
% scoring.d3_sigma = 0.15  (mal-core/src/mal_core/scoring/scorers/d3_eip.py:72)
\expandafter\def\csname cc@k@scoring.d3_sigma\endcsname{0.15}
% scoring.d4_collapse_max = 0.8  (mal-core/src/mal_core/scoring/scorers/d4_stability.py:70)
\expandafter\def\csname cc@k@scoring.d4_collapse_max\endcsname{0.8}
% scoring.d4_collapse_min = 0.1  (mal-core/src/mal_core/scoring/scorers/d4_stability.py:70)
\expandafter\def\csname cc@k@scoring.d4_collapse_min\endcsname{0.1}
% scoring.d4_explosion_max = 5.0  (mal-core/src/mal_core/scoring/scorers/d4_stability.py:77)
\expandafter\def\csname cc@k@scoring.d4_explosion_max\endcsname{5.0}
% scoring.d5_band_max = 0.85  (mal-core/src/mal_core/scoring/scorers/d5_morans.py:95)
\expandafter\def\csname cc@k@scoring.d5_band_max\endcsname{0.85}
% scoring.d5_band_min = 0.2  (mal-core/src/mal_core/scoring/scorers/d5_morans.py:95)
\expandafter\def\csname cc@k@scoring.d5_band_min\endcsname{0.2}
% scoring.d5_center = 0.5  (mal-core/src/mal_core/scoring/scorers/d5_morans.py:98)
\expandafter\def\csname cc@k@scoring.d5_center\endcsname{0.5}
% scoring.d5_sigma = 0.25  (mal-core/src/mal_core/scoring/scorers/d5_morans.py:98)
\expandafter\def\csname cc@k@scoring.d5_sigma\endcsname{0.25}
% scoring.d8_min_score = 0.4  (mal-core/src/mal_core/scoring/scorers/d8_coupling.py:35)
\expandafter\def\csname cc@k@scoring.d8_min_score\endcsname{0.4}
% scoring.d8_pass_threshold = 0.3  (mal-core/src/mal_core/scoring/scorers/d8_coupling.py:76)
\expandafter\def\csname cc@k@scoring.d8_pass_threshold\endcsname{0.3}
% scoring.d8_score_divisor = 0.5  (mal-core/src/mal_core/scoring/scorers/d8_coupling.py:70)
\expandafter\def\csname cc@k@scoring.d8_score_divisor\endcsname{0.5}
% scoring.gate_census_reference = 27000000  (mal-core/src/mal_core/scoring/scorers/gates.py:191)
\expandafter\def\csname cc@k@scoring.gate_census_reference\endcsname{27000000}
% scoring.gate_ratio_band_max = 1.2  (mal-core/src/mal_core/scoring/scorers/gates.py:214)
\expandafter\def\csname cc@k@scoring.gate_ratio_band_max\endcsname{1.2}
% scoring.gate_ratio_band_min = 0.8  (mal-core/src/mal_core/scoring/scorers/gates.py:214)
\expandafter\def\csname cc@k@scoring.gate_ratio_band_min\endcsname{0.8}
% scoring.gate_row_sum_threshold = 0.99  (mal-core/src/mal_core/scoring/scorers/gates.py:271)
\expandafter\def\csname cc@k@scoring.gate_row_sum_threshold\endcsname{0.99}
% scoring.gate_row_sum_tolerance = 0.01  (mal-core/src/mal_core/scoring/scorers/gates.py:260)
\expandafter\def\csname cc@k@scoring.gate_row_sum_tolerance\endcsname{0.01}
% scoring.gate_rural_class = 50  (mal-core/src/mal_core/scoring/scorers/gates.py:317)
\expandafter\def\csname cc@k@scoring.gate_rural_class\endcsname{50}
% scoring.gate_sample_files = 48  (mal-core/src/mal_core/scoring/scorers/gates.py:59)
\expandafter\def\csname cc@k@scoring.gate_sample_files\endcsname{48}
% scoring.gate_urban_class = 30  (mal-core/src/mal_core/scoring/scorers/gates.py:316)
\expandafter\def\csname cc@k@scoring.gate_urban_class\endcsname{30}
% scoring.gate_urban_ratio_max = 0.7  (mal-core/src/mal_core/scoring/scorers/gates.py:336)
\expandafter\def\csname cc@k@scoring.gate_urban_ratio_max\endcsname{0.7}
% scoring.gate_wall_time_limit_s = 30.0  (mal-core/src/mal_core/scoring/scorers/gates.py:158)
\expandafter\def\csname cc@k@scoring.gate_wall_time_limit_s\endcsname{30.0}
% scoring.regression_threshold = -0.05  (mal-core/src/mal_core/scoring/feedback.py:51)
\expandafter\def\csname cc@k@scoring.regression_threshold\endcsname{-0.05}
% scoring.verdict_borderline = 0.5  (mal-core/src/mal_core/scoring/runner.py:24)
\expandafter\def\csname cc@k@scoring.verdict_borderline\endcsname{0.5}
% scoring.verdict_regressed = 0.3  (mal-core/src/mal_core/scoring/runner.py:24)
\expandafter\def\csname cc@k@scoring.verdict_regressed\endcsname{0.3}
% scoring.verdict_viable = 0.7  (mal-core/src/mal_core/scoring/runner.py:24)
\expandafter\def\csname cc@k@scoring.verdict_viable\endcsname{0.7}
% training.ckpt_period = 10  (mal-core/src/mal_core/training/trainer.py:102)
\expandafter\def\csname cc@k@training.ckpt_period\endcsname{10}
% training.dice_count_threshold = 0.5  (mal-core/src/mal_core/training/model.py:110)
\expandafter\def\csname cc@k@training.dice_count_threshold\endcsname{0.5}
% training.dice_eps = 1e-06  (mal-core/src/mal_core/training/model.py:97)
\expandafter\def\csname cc@k@training.dice_eps\endcsname{1e-06}
% training.dice_tau = 0.5  (mal-core/src/mal_core/training/model.py:87)
\expandafter\def\csname cc@k@training.dice_tau\endcsname{0.5}
% training.dice_thresh = 0.5  (mal-core/src/mal_core/training/model.py:87)
\expandafter\def\csname cc@k@training.dice_thresh\endcsname{0.5}
% training.dice_weight = 0.5  (mal-core/src/mal_core/training/model.py:105)
\expandafter\def\csname cc@k@training.dice_weight\endcsname{0.5}
% training.env_channels = 4  (mal-core/src/mal_core/training/dataset.py:22)
\expandafter\def\csname cc@k@training.env_channels\endcsname{4}
% training.fallback_in_channels = 6  (mal-core/src/mal_core/training/wrapper.py:62)
\expandafter\def\csname cc@k@training.fallback_in_channels\endcsname{6}
% training.fallback_out_channels = 2  (mal-core/src/mal_core/training/wrapper.py:66)
\expandafter\def\csname cc@k@training.fallback_out_channels\endcsname{2}
% training.pad_multiple = 16  (mal-core/src/mal_core/training/wrapper.py:100)
\expandafter\def\csname cc@k@training.pad_multiple\endcsname{16}
% training.patch_size = 128  (mal-core/src/mal_core/training/dataset.py:19)
\expandafter\def\csname cc@k@training.patch_size\endcsname{128}
% training.sub_seed = 42  (mal-core/src/mal_core/training/dataset.py:210)
\expandafter\def\csname cc@k@training.sub_seed\endcsname{42}
% training.transmission_channels = 4  (mal-core/src/mal_core/training/dataset.py:21)
\expandafter\def\csname cc@k@training.transmission_channels\endcsname{4}
% training.unet_channels = 32, 64, 128, 256  (mal-core/src/mal_core/training/model.py:46)
\expandafter\def\csname cc@k@training.unet_channels\endcsname{32, 64, 128, 256}
% training.vector_channels = 2  (mal-core/src/mal_core/training/dataset.py:20)
\expandafter\def\csname cc@k@training.vector_channels\endcsname{2}


% ---------------------------------------------------------------------------
% Values that do not live in mal-core/config/model_constants.yaml.  Every one
% of them is a literal in the source of the training package; each is listed,
% with its file and line, in the evidence appendix ("Code literals used in
% this document").  tools/gen_concepts_values.py reads them out of the code
% text and emits them as \codeconst{...} macros in
% docs/concepts/generated/code-macros.tex, so that a value the code changes
% makes this document fail to build instead of making it lie.  The macros
% below are short aliases for those generated macros; the body of the
% document never contains a hand-typed number.
% ---------------------------------------------------------------------------
\newcommand{\PatchSize}{\codeconst{training.patch_size}}
\newcommand{\EnvChannels}{\codeconst{training.env_channels}}
\newcommand{\VectorChannels}{\codeconst{training.vector_channels}}
\newcommand{\TransChannels}{\codeconst{training.transmission_channels}}
\newcommand{\DiceWeight}{\codeconst{training.dice_weight}}
\newcommand{\DiceTau}{\codeconst{training.dice_tau}}
\newcommand{\DiceThresh}{\codeconst{training.dice_thresh}}
\newcommand{\DiceEps}{\ensuremath{\text{\codeconst{training.dice_eps}}}}
\newcommand{\SubSeed}{\codeconst{training.sub_seed}}
\newcommand{\CkptPeriod}{\codeconst{training.ckpt_period}}
\newcommand{\PoolMult}{\codeconst{training.pad_multiple}}
\newcommand{\StateContract}{\mskey{contracts.cogs_state}}
\newcommand{\EnvContract}{\mskey{contracts.env_nc}}

\title{The \emph{training} stage of the Centinela:\\
a learned one-step surrogate for the agent-based model}
\author{Conceptual document, \texttt{docs/concepts/training/}}
\date{}

\begin{document}
\maketitle

\begin{abstract}
The \emph{training} stage fits a small convolutional network --- a U-Net --- to
imitate one simulated day of the agent-based model (ABM). It reads the daily
state snapshots that an ABM run leaves on disk, tiles them, pairs consecutive
days, splits the tiles spatially, and optimises the network so that its
prediction of tomorrow's state matches what the ABM actually produced
tomorrow. The fitted weights are written out as checkpoints that a later
stage (\emph{prediction}) loads --- but only when the operator names the
model explicitly --- in order to produce a risk map without running the C++
ABM at all.

This document explains what the stage is for, what it consumes, what it
produces, in what order it does its work, how the data is assembled and split,
what the network is and what it is optimising, how the ``best'' model is
chosen, what a checkpoint contains and how it is read back, and where the
seam between training and prediction is thin. It is written for a reader who
has never opened the code, and every factual claim resolves to a file and a
line.
\end{abstract}

\tableofcontents

% ===========================================================================
\section{Purpose, and what this stage is not}
\label{sec:purpose}

\subsection{The question the stage exists to answer}

Running the full agent-based model is expensive. The ABM is a C++ binary that
must be executed day by day for every area of interest, every year, every
scenario; a decision-support system that has to answer ``where is transmission
risk concentrated in this district next month'' cannot afford to re-run it for
every question. The \emph{training} stage therefore exists to answer one
question:

\begin{quote}
\emph{Can a small learned network reproduce the ABM's own one-day state
transition closely enough that the prediction stage can use it in place of
re-running the ABM?}
\end{quote}

The stage's own command-line help states the intent in exactly these terms:
it trains a U-Net ``to predict ABM outcomes from environmental tensors,
enabling fast inference without full ABM runs''
(\codex{mal-core/\allowbreak{}src/\allowbreak{}mal\_\allowbreak{}core/\allowbreak{}cli.py:784-785}). The prediction stage can close the
loop, but it does not do so by itself: its \codex{--model} option defaults to
\codex{"dummy"}
(\codex{mal-core/\allowbreak{}src/\allowbreak{}mal\_\allowbreak{}core/\allowbreak{}cli.py:112}), and the registry answers that name with a
zero-emitting \codex{DummyModel}
(\codex{mal-core/\allowbreak{}src/\allowbreak{}mal\_\allowbreak{}core/\allowbreak{}prediction/\allowbreak{}registry.py:41-46,96-97}). A trained
checkpoint is therefore \emph{never picked up automatically}: the surrogate
is applied only when the operator passes the model's registry key to
\codex{--model}, in which case the registry resolves the manifest and builds
a wrapper around the declared checkpoint
(\codex{mal-core/\allowbreak{}src/\allowbreak{}mal\_\allowbreak{}core/\allowbreak{}prediction/\allowbreak{}registry.py:96-118}).

\subsection{What would count as the question having been answered}

The question is answered in the operational sense when a checkpoint exists on
disk that the prediction stage can find, load and apply to a state snapshot.
Concretely, the training stage has done its job when the output directory
contains a manifest file named \codex{model.yaml} naming a checkpoint file that
exists next to it (\codex{mal-core/\allowbreak{}src/\allowbreak{}mal\_\allowbreak{}core/\allowbreak{}training/\allowbreak{}trainer.py:88,93,105-106}),
the prediction stage's registry can scan that directory, resolve the
manifest, and construct the network (\codex{mal-core/\allowbreak{}src/\allowbreak{}mal\_\allowbreak{}core/\allowbreak{}prediction/\allowbreak{}registry.py:61-74,96-118}),
and the operator names that model on the prediction command line --- the
\codex{--model} default of \codex{"dummy"} would otherwise select the
zero-emitting dummy
(\codex{mal-core/\allowbreak{}src/\allowbreak{}mal\_\allowbreak{}core/\allowbreak{}cli.py:112};
\codex{mal-core/\allowbreak{}src/\allowbreak{}mal\_\allowbreak{}core/\allowbreak{}prediction/\allowbreak{}registry.py:96-97}).

Notice what this criterion does \emph{not} include. It does not include any
accuracy floor, any comparison against the ABM, any tolerance band, any
acceptance scorecard. \emph{The stage answers the question when it has produced
a loadable artefact; whether that artefact is faithful to the ABM is not
checked anywhere.} Section~\ref{sec:absences} returns to this.

\subsection{Negative purpose: what the surrogate is not}

The surrogate introduced in Section~\ref{sec:surrogate} is a very particular
object, and the following are \emph{not} among its properties. They are stated
here, at the point where the reader first meets the surrogate, because a reader
who knows them will read the rest of this document correctly.

\begin{itemize}
\item \textbf{It is not a measurement.} The network's output is a set of
      learned parameters. It has no physical unit of its own; the units it
      carries are borrowed from whatever the ABM happened to write into the
      training rasters, and even those are inherited by imitation, not by
      construction. Nothing in the training path measures anything.
\item \textbf{It is not a model of the ABM.} It is a model of \emph{one ABM
      step}, in one area, under the conditions the supplied rollouts
      exhibited. It has no representation of the mechanisms the ABM
      implements --- no gonotrophic cycle, no host-seeking, no larval
      competition, no dispersal kernel, no human immunity. Those mechanisms
      survive only as statistical regularities in the training pairs.
\item \textbf{It is not a projection.} Nothing in the code rolls the surrogate
      forward over a horizon. The prediction stage calls it exactly once per
      invocation and aggregates the single answer into a risk map
      (\codex{mal-core/\allowbreak{}src/\allowbreak{}mal\_\allowbreak{}core/\allowbreak{}prediction/\allowbreak{}predictor.py:89-92}). The
      surrogate cannot produce a multi-day trajectory, a prevalence curve, or
      a scenario comparison.
\item \textbf{It is not a generalisation beyond the rollouts it saw.} It is
      trained only on the ABM outputs present in the run directory it was
      given. What it can be expected to represent \emph{outside} those
      rollouts --- a different intervention regime, a different area of
      interest, a different year, a newer version of the ABM --- is a
      statistical expectation: nothing in this repository establishes it,
      and no artefact of this stage records it (see the evidence appendix).
\item \textbf{It is not a calibrated instrument.} There is no scorer, no
      tolerance, no floor and no scorecard for this stage (Section
      \ref{sec:absences}). The word ``best'' in the training loop means
      ``the largest validation Dice seen so far, with ties overwriting'',
      nothing more.
\end{itemize}

\subsection{Which source wins when they disagree}

This document draws its values from three sources, and they do not carry equal
authority. \textbf{Values} --- the numbers that enter a computation --- are
governed by the shared configuration file
(\codex{mal-core/\allowbreak{}config/\allowbreak{}model\_\allowbreak{}constants.yaml}), and that file wins.
\textbf{Contract values and names} --- version strings, band names, artefact
names --- are governed by the generated contract table, and that table wins.
\textbf{Semantics} --- what the numbers mean, and how the stage behaves --- are
governed by this document, because this document owns the training stage's
prose, and this document wins. When two of these sources contradict each
other, the one with the lower authority is the bug: a value that disagrees
with the configuration file is an error in the configuration file or in the
code that reads it, a contract name that disagrees with the generated table
is an error in one of them, and a sentence here that disagrees with the code
is an error in the sentence. The rule is stated in this document rather than
on a shared page because the training stage's authority belongs to this
document, and no other document speaks for it.

% ===========================================================================
\section{The surrogate, introduced}

\label{sec:surrogate}

\subsection{The mapping}

The surrogate is a feed-forward convolutional network, a \emph{U-Net}
(a U-shaped encoder--decoder network with skip connections), implemented as
the class \codex{UNet} (\codex{mal-core/\allowbreak{}src/\allowbreak{}mal\_\allowbreak{}core/\allowbreak{}training/\allowbreak{}model.py:28-46}).
It learns a one-step state transition:

\begin{equation}
\widehat{S_{t+1}} \;=\; f_\theta\!\left(S_t \,\Vert\, T_t \,\Vert\, E\right),
\label{eq:mapping}
\end{equation}

where \(S_t\) is the ABM's vector state at day \(t\), \(T_t\) is the optional
transmission state at day \(t\), \(E\) is the environmental context, and
\(\Vert\) denotes channel-wise concatenation. The target of the fit is the
ABM's own state at day \(t+1\), \(S_{t+1}\) (and \(T_{t+1}\) where present).

The \emph{temporal horizon of one training example is exactly one simulated
day}. Pairs are formed from consecutive snapshots present in the run
directory: the code sorts the day numbers of a rollout and pairs
\((d_i, d_{i+1})\) for each adjacent pair in that sorted list
(\codex{mal-core/\allowbreak{}src/\allowbreak{}mal\_\allowbreak{}core/\allowbreak{}training/\allowbreak{}dataset.py:135-146}). Two consequences
follow and both matter. First, the horizon is a day, so a surrogate that has
been trained cannot be asked for a week's transition. Second, the pairing is
by \emph{adjacency in the sorted list of present days}, not by adjacency in
integers: if a run directory holds days \(1,2,9,10\), the pair \((2,9)\) is
formed, and it spans an interval of seven days even though the network is
being taught a one-day transition.

\subsection{Channel composition}

The channel counts are \emph{not} fixed constants in the training path; they
are read off the data at run time. The trainer takes a single sample from the
training dataset and reads the channel counts off the sample's two axes
(\codex{mal-core/\allowbreak{}src/\allowbreak{}mal\_\allowbreak{}core/\allowbreak{}training/\allowbreak{}trainer.py:50-52}), then builds the
network with those counts (\codex{mal-core/\allowbreak{}src/\allowbreak{}mal\_\allowbreak{}core/\allowbreak{}training/\allowbreak{}trainer.py:54}).
The composition is:

\begin{center}
\begin{tabular}{@{}lll@{}}
\toprule
Tensor & Channels & Source \\
\midrule
Input \(x\) & \(\VectorChannels\) (state) \(\oplus\) \(\TransChannels\) (transmission, optional) \(\oplus\) \EnvChannels{} (environment)
  & \codex{dataset.py:245-252} \\
Target \(y\) & \(\VectorChannels\) (state) \(\oplus\) \(\TransChannels\) (transmission, optional)
  & \codex{dataset.py:242-243} \\
\bottomrule
\end{tabular}
\end{center}

With transmission enabled --- the default --- this is a ten-channel input and a
six-channel target; without it, six in and two out. The module's own
docstring still advertises the older six-in/two-out shape
(\codex{mal-core/\allowbreak{}src/\allowbreak{}mal\_\allowbreak{}core/\allowbreak{}training/\allowbreak{}model.py:1-5}), and the constructor's
default arguments are still the \codex{in\_\allowbreak{}channels} and
\codex{out\_\allowbreak{}channels} values recorded in the generated parameter
table of Section~\ref{sec:architecture}
(\codex{mal-core/\allowbreak{}src/\allowbreak{}mal\_\allowbreak{}core/\allowbreak{}training/\allowbreak{}model.py:29-30}); the training loop
overrides both from the data, so the docstring is stale but harmless in the
training path. It is \emph{not} harmless in the standalone validation script,
which constructs the network with those defaults and would fail to load a
transmission-enabled checkpoint (Section~\ref{sec:absences}).

Two module-level constants, \codex{VECTOR\_\allowbreak{}CHANNELS} and
\codex{TRANSMISSION\_\allowbreak{}CHANNELS}, are declared in the dataset module
(\codex{mal-core/\allowbreak{}src/\allowbreak{}mal\_\allowbreak{}core/\allowbreak{}training/\allowbreak{}dataset.py:20-21}) but never read
anywhere; only \codex{ENV\_\allowbreak{}CHANNELS} is used
(\codex{mal-core/\allowbreak{}src/\allowbreak{}mal\_\allowbreak{}core/\allowbreak{}training/\allowbreak{}dataset.py:22}). The state and
transmission band counts are in fact taken from the rasters themselves
(\codex{mal-core/\allowbreak{}src/\allowbreak{}mal\_\allowbreak{}core/\allowbreak{}training/\allowbreak{}dataset.py:179,229-230}). A consequence
worth stating: a run directory whose rasters carry an unexpected number of
bands produces a model with that number of channels, and nothing complains.

\subsection{What the surrogate cannot represent}

The limitations are structural consequences of the three facts above, and the
reader should hold them in mind while reading Sections
\ref{sec:data}--\ref{sec:selection}.

\begin{enumerate}
\item \textbf{It is purely local.} The network is convolutional with a
      receptive field of finite extent; a cell's prediction depends on its
      neighbourhood, not on the whole raster. It cannot express a
      population-level constraint or a basin-wide balance.
\item \textbf{It has no memory.} The input is a single state snapshot plus a
      single static environmental stack. There is no recurrent state, no
      stack of previous days, and no daily rainfall history. A rainfall event
      on day \(t\) cannot influence the prediction for day \(t+1\), because
      the environmental planes the network sees are the same for every
      training pair (Section~\ref{sec:env}).
\item \textbf{It is trained on one regime only.} It inherits the ABM's
      behaviour for the seeds, years, areas and interventions present in the
      supplied run directory, and nothing else.
\item \textbf{It is bounded only by what the loss encourages.} The output
      head is a bare convolution with no activation
      (\codex{mal-core/\allowbreak{}src/\allowbreak{}mal\_\allowbreak{}core/\allowbreak{}training/\allowbreak{}model.py:46}), so a predicted
      adult count may come out negative or fractional, and nothing clamps it
      in training or at inference
      (\codex{mal-core/\allowbreak{}src/\allowbreak{}mal\_\allowbreak{}core/\allowbreak{}training/\allowbreak{}wrapper.py:106}). How the two
      loss terms actually traded off during the fit is not observable: the
      code records no per-epoch loss beyond the single printed line
      (\codex{mal-core/\allowbreak{}src/\allowbreak{}mal\_\allowbreak{}core/\allowbreak{}training/\allowbreak{}trainer.py:84}).
\end{enumerate}

% ===========================================================================
\section{Notation}

The table below fixes the symbols used in this document. The parenthesis gives
the code name, so that a reader can hold the description and the source in one
head.

\begin{center}
\begin{longtable}{@{}lp{3.4cm}p{9.0cm}@{}}
\toprule
Symbol & Code name & Meaning \\
\midrule
\endhead
\(S_t\) & state raster, bands 1--2 & The ABM's vector state at day \(t\): two bands of raw per-cell adult counts, \codex{adult\_\allowbreak{}patch\_\allowbreak{}count} and \codex{adult\_\allowbreak{}cell\_\allowbreak{}count} (\codex{mal-core/\allowbreak{}src/\allowbreak{}mal\_\allowbreak{}core/\allowbreak{}abm/\allowbreak{}include/\allowbreak{}mal\_\allowbreak{}abm\_\allowbreak{}fast/\allowbreak{}wire.hpp:133}), under state contract \StateContract. \\
\(T_t\) & transmission raster, bands 1--4 & The ABM's transmission state at day \(t\): four bands, named \codex{human\_\allowbreak{}prevalence}, \codex{human\_\allowbreak{}incidence}, \codex{infectious\_\allowbreak{}vector\_\allowbreak{}pressure} and \codex{active\_\allowbreak{}focus} by the transmission output writer (\codex{mal-core/\allowbreak{}src/\allowbreak{}mal\_\allowbreak{}core/\allowbreak{}abm/\allowbreak{}include/\allowbreak{}mal\_\allowbreak{}abm\_\allowbreak{}fast/\allowbreak{}transmission\_\allowbreak{}output.hpp:21-26}). The fourth is \emph{not} the ``transmission focus'' band of the prediction risk map, which is a different artefact (\codex{mal-core/\allowbreak{}src/\allowbreak{}mal\_\allowbreak{}core/\allowbreak{}prediction/\allowbreak{}predictor.py:39}). \\
\(E\) & environmental stack & Up to four planes, read once from a NetCDF or GeoTIFF supplied by the operator (\codex{mal-core/\allowbreak{}src/\allowbreak{}mal\_\allowbreak{}core/\allowbreak{}training/\allowbreak{}dataset.py:66-88}); with the real environmental NetCDF the read fails and every plane is zero (Section~\ref{sec:env}). \\
\(f_\theta\) & \codex{UNet} & The surrogate network (\codex{mal-core/\allowbreak{}src/\allowbreak{}mal\_\allowbreak{}core/\allowbreak{}training/\allowbreak{}model.py:28-60}). \\
\(\theta\) & \codex{state\_\allowbreak{}dict} & The learned parameters, as written to a checkpoint. \\
\(P\) & \codex{patch\_\allowbreak{}size} & The side of a training tile, in raster cells (\codex{mal-core/\allowbreak{}src/\allowbreak{}mal\_\allowbreak{}core/\allowbreak{}training/\allowbreak{}dataset.py:19}). \\
\(c_{\mathrm{mid}}\) & \codex{mid\_\allowbreak{}col} & The column at which the raster is split between training and validation tiles (\codex{mal-core/\allowbreak{}src/\allowbreak{}mal\_\allowbreak{}core/\allowbreak{}training/\allowbreak{}dataset.py:131}). \\
\(D_\mathrm{train}, D_\mathrm{val}\) & \codex{RolloutDataset} & The two splits, each a list of tile records (\codex{mal-core/\allowbreak{}src/\allowbreak{}mal\_\allowbreak{}core/\allowbreak{}training/\allowbreak{}dataset.py:99-170}). \\
\(\mathrm{MSE}\) & \codex{F.mse\_\allowbreak{}loss} & The mean squared error over \emph{all} channels (\codex{mal-core/\allowbreak{}src/\allowbreak{}mal\_\allowbreak{}core/\allowbreak{}training/\allowbreak{}model.py:79}). \\
\(\mathrm{Dice}_{\mathrm{soft}}\) & \codex{soft\_\allowbreak{}dice\_\allowbreak{}loss} & The sigmoid-smoothed Dice on the first channel (\codex{mal-core/\allowbreak{}src/\allowbreak{}mal\_\allowbreak{}core/\allowbreak{}training/\allowbreak{}model.py:63-74}). \\
\(\mathrm{Dice}_{\mathrm{hard}}\) & \codex{eval\_\allowbreak{}dice} & The hard Dice on the first channel at a fixed count threshold (\codex{mal-core/\allowbreak{}src/\allowbreak{}mal\_\allowbreak{}core/\allowbreak{}training/\allowbreak{}model.py:86-95}). \\
\bottomrule
\end{longtable}
\end{center}

% ===========================================================================
\section{What the stage consumes}
\label{sec:consumes}

\subsection{The ABM's daily snapshots}

The training data are the daily snapshots an ABM run writes into its output
directory. Two conventions of the ABM's file naming are load-bearing here,
because the training stage discovers its data by globbing rather than by
reading a manifest:

\begin{itemize}
\item The per-day state raster is \codex{<stem>\_\allowbreak{}dayNNN.tif}, with the day
      number zero-padded to three digits
      (\codex{mal-core/\allowbreak{}src/\allowbreak{}mal\_\allowbreak{}core/\allowbreak{}abm/\allowbreak{}src/\allowbreak{}main.cpp:173-185}).
\item The per-day transmission raster follows the same
      \codex{<stem>\_\allowbreak{}dayNNN.tif} pattern
      (\codex{mal-core/\allowbreak{}src/\allowbreak{}mal\_\allowbreak{}core/\allowbreak{}abm/\allowbreak{}src/\allowbreak{}main.cpp:218-230}); the two
      families are distinguished only by the substring
      \codex{transmission} appearing in the file name.
\end{itemize}

The state rasters carry \emph{raw per-cell adult counts}, not densities
normalised by any capacity constant. Under state contract \StateContract the
retired \codex{K\_\allowbreak{}MAX} normalisation is absent from the whole output path
(\codex{mal-core/\allowbreak{}src/\allowbreak{}mal\_\allowbreak{}core/\allowbreak{}abm/\allowbreak{}src/\allowbreak{}output\_\allowbreak{}contract.cpp:272}); the surrogate
is therefore fitted to the same quantity the ABM emits, and no normalisation
step exists anywhere in the training package.

\subsection{The environmental context}

The environmental context is not discovered; the operator names it. The
command-line option \codex{--env-path} accepts a NetCDF or a GeoTIFF
(\codex{mal-core/\allowbreak{}src/\allowbreak{}mal\_\allowbreak{}core/\allowbreak{}cli.py:780}), and the dataset reads it once, at
construction (\codex{mal-core/\allowbreak{}src/\allowbreak{}mal\_\allowbreak{}core/\allowbreak{}training/\allowbreak{}dataset.py:66-88}). Its
contents and its ordering are treated in Section~\ref{sec:env}.

\subsection{What the run directory must contain for the stage to run at all}

The stage requires, at minimum, a run directory containing files whose names
carry the substring \codex{\_\allowbreak{}day} and one of the extensions \codex{.tif} or
\codex{.npy}. If none is found, the dataset raises
\codex{FileNotFoundError} (\codex{mal-core/\allowbreak{}src/\allowbreak{}mal\_\allowbreak{}core/\allowbreak{}training/\allowbreak{}dataset.py:99-101}).
The generated contract table below places this stage among the consumers of
the ABM's state output.

\begin{landscape}
{\tiny\setlength{\tabcolsep}{3pt}\setlength{\LTleft}{0pt}\setlength{\LTright}{0pt}%
% Generated by tools/gen_concepts_values.py -- do not edit.
\paragraph{The output contracts}
\begin{longtable}{@{}llllll@{}}
\toprule
Artefact & Producer & Consumers & Version & Bands / variables & Sentinel \\
\midrule
\endhead
\code{env} & \code{mal\_core.ingest.env.build\_env\_tensor} & abm, prediction, scoring & "3.0" & rainfall, air\_temp\_c, jrc\_persistent, lc\_built, lc\_wetland, lc\_open\_plus\_water, lc\_crop, land\_frac (optional), salinity\_ppt (optional) & none \\
\code{habitat} & \code{mal\_core.ingest.env.build\_env\_tensor} & abm &  & hab\_type, habitat\_potential, row, col & none \\
\code{state\_cog} & \code{mal\_abm\_fast::write\_state\_cog} & scoring, prediction, training & "4.0" & adult\_patch\_count, adult\_cell\_count & -9999.0 \\
\code{state\_sidecar} & \code{mal\_abm\_fast::write\_state\_sidecar} & scoring & "4.0" &  &  \\
\code{manifest} & \code{mal\_core.download.manifest.update\_dataset} & download, ingest, abm &  &  &  \\
\code{host\_static} & \code{mal\_core.ingest.hosts.build\_host\_dataset} & abm &  & humans, livestock, urban, buildings &  \\
\code{mobility\_day} & \code{mal\_core.ingest.mobility.build\_mobility\_dataset} & abm &  &  &  \\
\bottomrule
\end{longtable}

\paragraph{The state sidecar's keys, in emission order}
\begin{longtable}{@{}l@{}}
\toprule
Key \\
\midrule
\endhead
\code{crs} \\
\code{transform} \\
\code{aoi\_slug} \\
\code{scale} \\
\code{year} \\
\code{month} \\
\code{seed} \\
\code{n\_rollouts} \\
\code{rollout\_index} \\
\code{generator\_version} \\
\code{abm\_params\_hash} \\
\code{contract\_version} \\
\code{band\_names} \\
\code{nodata} \\
\code{shape} \\
\bottomrule
\end{longtable}
}
\end{landscape}

% ===========================================================================
\section{What the stage produces, and who consumes it}
\label{sec:produces}

\subsection{The four artefacts}

The training stage writes into its output directory:

\begin{center}
\begin{longtable}{@{}lp{3.2cm}p{3.7cm}p{3.7cm}@{}}
\toprule
Artefact & Written at & Contents & Consumed by \\
\midrule
\endhead
\codex{best\_\allowbreak{}model.pt} & \codex{trainer.py:88} & Bare \codex{state\_\allowbreak{}dict} of the best-scoring weights & \codex{mal-core/\allowbreak{}src/\allowbreak{}mal\_\allowbreak{}core/\allowbreak{}prediction/\allowbreak{}registry.py:110-114} (as the manifest's declared checkpoint) \\
\codex{model\_\allowbreak{}epoch\_\allowbreak{}N.pt} & \codex{trainer.py:90-91} & Bare \codex{state\_\allowbreak{}dict}, written every \CkptPeriod{}th epoch & Nothing in the tree; a manual inspection point \\
\codex{final\_\allowbreak{}model.pt} & \codex{trainer.py:93} & Bare \codex{state\_\allowbreak{}dict} of the last epoch's weights & Nothing in the tree \\
\codex{model.yaml} & \codex{trainer.py:96-106} & A small YAML manifest & The prediction stage's model registry (\codex{registry.py:61-74}) \\
\bottomrule
\end{longtable}
\end{center}

\subsection{The manifest}

The manifest carries seven keys and no more
(\codex{mal-core/\allowbreak{}src/\allowbreak{}mal\_\allowbreak{}core/\allowbreak{}training/\allowbreak{}trainer.py:96-105}):

\begin{center}
\begin{tabular}{@{}lp{3.2cm}p{8.2cm}@{}}
\toprule
Key & Value & Meaning \\
\midrule
\codex{name} & the output directory's own name & The model is named after its directory, so two training runs into differently named directories yield two differently named models. \\
\codex{version} & \codex{\codeconst{prediction.manifest_version}} & The model's own version; the registry keys models as \codex{name@version}. \\
\codex{contract\_\allowbreak{}version} & \codex{\codeconst{prediction.contract_version_manifest}} & The \emph{prediction} contract version, written as a hard-coded literal by the trainer (\codex{mal-core/\allowbreak{}src/\allowbreak{}mal\_\allowbreak{}core/\allowbreak{}training/\allowbreak{}trainer.py:99}); it is not the ABM state contract, and it is not read from the registry --- the registry's own default of \codex{\codeconst{prediction.contract_version_manifest}} is only a fallback for manifests that omit the key (\codex{mal-core/\allowbreak{}src/\allowbreak{}mal\_\allowbreak{}core/\allowbreak{}prediction/\allowbreak{}registry.py:24-31}). \\
\codex{in\_\allowbreak{}channels} & read off the data & Number of input channels the network was built with. \\
\codex{out\_\allowbreak{}channels} & read off the data & Number of output channels. \\
\codex{checkpoint} & \codex{best\_\allowbreak{}model.pt} & Which of the three checkpoints the registry should load. \\
\codex{description} & free text & Contains the run directory path and the channel counts, and nothing else machine-readable. \\
\bottomrule
\end{tabular}
\end{center}

Two facts a reader should notice by their absence. The manifest does not
record the ABM state-contract version under which the training rasters were
written --- so a checkpoint fitted to contract-\StateContract{} raw counts is
indistinguishable from one fitted to any other contract --- and it records no
provenance of the training data beyond a path embedded in a sentence: not the
area of interest, not the seed set, not the number of days, not the patch
size, not the split rule, not the date of training.

\subsection{The consumption contract, stated exactly}

The prediction stage consumes the checkpoint through the manifest, and the
contract is thin: the manifest must sit in the same directory as the
checkpoint it names, the checkpoint must exist, and the channel counts in the
manifest must match the weights in the checkpoint. Given those, the wrapper
rebuilds the network and loads the weights strictly
(\codex{mal-core/\allowbreak{}src/\allowbreak{}mal\_\allowbreak{}core/\allowbreak{}training/\allowbreak{}wrapper.py:64-67}); a mismatch raises at
construction. The registry additionally requires the checkpoint to exist on
disk, raising \codex{FileNotFoundError} otherwise
(\codex{mal-core/\allowbreak{}src/\allowbreak{}mal\_\allowbreak{}core/\allowbreak{}prediction/\allowbreak{}registry.py:106-109}).

% ===========================================================================
\section{Order of operations, and its consequences}

\subsection{The order}

\begin{longtable}{@{}cp{2.4cm}p{8.6cm}p{1.9cm}@{}}
\toprule
Step & Actor & Order & Mode \\
\midrule
\endhead
1 & \codex{train\_\allowbreak{}unet} & Create the output directory (\codex{trainer.py:26-27}) & state \\
2 & \codex{train\_\allowbreak{}unet} & Resolve the device: CUDA, else Apple MPS, else CPU (\codex{trainer.py:29-35}) & choice \\
3 & dataset & Read the environmental stack, once (\codex{dataset.py:61,66-88}) & input \\
4 & dataset & Discover, group, pair, tile, split, subsample (\codex{dataset.py:90-170}) & derived \\
5 & \codex{train\_\allowbreak{}unet} & Build the two data loaders (\codex{trainer.py:37-45}) & derived \\
6 & \codex{train\_\allowbreak{}unet} & Raise if the training split is empty (\codex{trainer.py:47-48}) & check \\
7 & \codex{train\_\allowbreak{}unet} & Read the channel counts off one training sample (\codex{trainer.py:50-52}) & derived \\
8 & \codex{train\_\allowbreak{}unet} & Build the network and the optimiser (\codex{trainer.py:54-55}) & state \\
9 & \codex{train\_\allowbreak{}unet} & Epoch loop: train, evaluate, select, checkpoint (\codex{trainer.py:57-91}) & state \\
10 & \codex{train\_\allowbreak{}unet} & Write the final checkpoint and the manifest (\codex{trainer.py:93-106}) & output \\
\bottomrule
\end{longtable}

\subsection{Consequences of that order}

\begin{itemize}
\item \textbf{The environment is read before the data is known to exist.}
      Step 3 precedes step 4, so a malformed environment file is consumed
      and discarded before the stage discovers that the run directory is
      empty. The environmental failure is silent
      (\codex{dataset.py:87-88}); the empty-directory failure is loud
      (\codex{dataset.py:99-101}). The two errors do not look alike to the
      operator, although the first is the more consequential.
\item \textbf{The split is decided before the subsample is drawn.} Steps 4
      contains both the mid-column split and the random subsample, in that
      order (\codex{dataset.py:148-170}). The subsample is therefore drawn
      independently within each split, each from its own patch list, and each
      from a generator re-seeded to the same fixed value
      (\codex{dataset.py:167}). The subsets are reproducible across runs on
      the same directory, but they are not the same subset of tiles --- they
      are different subsets of different populations.
\item \textbf{The channel counts are frozen from one sample.} Step 7 reads
      the widths of a single element of the training set and step 8 builds
      the network for those widths for the whole run. If the first training
      tile happened to lack transmission channels while later tiles had them,
      the network would be built too narrow and the wider tiles would fail at
      the concatenation. In practice the series are generated together, so
      the widths are uniform; the code does not check that they are.
\item \textbf{Selection happens inside the epoch loop, not after it.} Step 9
      both optimises and decides which weights to keep, so the ``best''
      checkpoint is a running maximum over epochs rather than a model chosen
      from a completed run (Section~\ref{sec:selection}).
\item \textbf{The manifest is written last.} Step 10 runs after the final
      checkpoint, so a run interrupted before step 10 leaves weights with no
      manifest, and the prediction stage's registry, which discovers models
      by manifest, will not see them
      (\codex{mal-core/\allowbreak{}src/\allowbreak{}mal\_\allowbreak{}core/\allowbreak{}prediction/\allowbreak{}registry.py:61-74}).
\end{itemize}

% ===========================================================================
\section{Assembling the training data}

\label{sec:data}

\subsection{Discovery}

The dataset globs the run directory for all \codex{*.tif} files, keeps those
whose names contain \codex{\_\allowbreak{}day} and do not contain \codex{transmission}, and
separately globs for \codex{*\_\allowbreak{}day*.npy}
(\codex{mal-core/\allowbreak{}src/\allowbreak{}mal\_\allowbreak{}core/\allowbreak{}training/\allowbreak{}dataset.py:90-97}). GeoTIFF is preferred
when the two families are of equal size or the GeoTIFF family is larger; the
choice is global for the whole dataset, not per file. The two families are
thus interchangeable only in the sense that the newer GeoTIFF path wins.

The discovery is looser than the naming convention the ABM follows
(Section~\ref{sec:consumes}): any file with \codex{\_\allowbreak{}day} in its name and no
\codex{transmission} substring is accepted as a state snapshot, and bands are
read positionally (\codex{dataset.py:178-179}), so the band \emph{order} is
assumed rather than checked against the band descriptions the ABM writes
(\codex{mal-core/\allowbreak{}src/\allowbreak{}mal\_\allowbreak{}core/\allowbreak{}abm/\allowbreak{}src/\allowbreak{}output\_\allowbreak{}contract.cpp:202,223}).

\subsection{Seeds, days, and pairing}

The seed and the day are recovered from the file name by regular expression:
\codex{seed(\d+)} and \codex{\_\allowbreak{}day(\d+)}
(\codex{mal-core/\allowbreak{}src/\allowbreak{}mal\_\allowbreak{}core/\allowbreak{}training/\allowbreak{}dataset.py:25-36}). A file whose name
carries neither pattern is assigned seed one and day zero --- a fallback, not
a parse. Snapshots are then grouped into a per-seed dictionary indexed by day
(\codex{dataset.py:102-110}), and the training pairs are the adjacent entries
of the sorted day list (\codex{dataset.py:135-146}).

The transmission rasters are indexed the same way
(\codex{dataset.py:112-118}) and are attached to a pair only when \emph{both}
the current and the next day's transmission raster exist for that pair
(\codex{dataset.py:143-146}); the concatenation itself happens at
\codex{dataset.py:242-243}. A pair without transmission is therefore a
state-only pair. Because the ABM writes the two families together, the
width of the dataset is in practice uniform; the code does not enforce it.

\subsection{Tiling}

The raster is cut into square tiles of side \(P = \PatchSize\)
(\codex{mal-core/\allowbreak{}src/\allowbreak{}mal\_\allowbreak{}core/\allowbreak{}training/\allowbreak{}dataset.py:19}). The tile grid counts
are \(\lfloor H/P \rfloor\) rows and \(\lfloor W/P \rfloor\) columns
(\codex{dataset.py:129-130}), so a raster whose height or width is not an
exact multiple of \(P\) has its remainder \emph{not represented in training
at all}: those cells are never cropped, never seen. The cropping and the
zero-padding of the edge tile are implemented in the item accessor
(\codex{dataset.py:210-225}), but for a grid computed by integer division the
edge tile is already full-size, so the padding branch is a safety net rather
than a routine path.

\subsection{No-data and dtype}

Rasters are read as \codex{float32}, and cells whose value equals the file's
declared no-data sentinel are set to zero
(\codex{mal-core/\allowbreak{}src/\allowbreak{}mal\_\allowbreak{}core/\allowbreak{}training/\allowbreak{}dataset.py:181-182}). The \codex{.npy}
fallback is cast to \codex{float32} and given a channel axis if it is
two-dimensional (\codex{dataset.py:184-186}). The consequence is worth naming:
\textbf{no-data becomes zero, which is the same value as ``no adults''}. The
surrogate cannot distinguish an unobserved cell from an empty one, and no
mask is carried through to the target.

\subsection{Subsampling}

When the subsample fraction is below one, a reproducible random subset of the
tiles is kept: the generator is seeded with a fixed value
(\codex{dataset.py:167}), the number kept is
\(\max(1, \lfloor \lvert \text{tiles} \rvert \times \text{fraction} \rfloor)\),
and the indices are drawn without replacement
(\codex{dataset.py:168-169}). Because the seed is fixed, the same subset is
drawn on every run over the same run directory.

% ===========================================================================
\section{The environmental stack, and its problems}

\label{sec:env}

\subsection{What is read}

The environmental context is read once, at dataset construction, into a cache
(\codex{mal-core/\allowbreak{}src/\allowbreak{}mal\_\allowbreak{}core/\allowbreak{}training/\allowbreak{}dataset.py:59,66-88}).

\begin{itemize}
\item From a \emph{NetCDF}, the variables are drawn in the order
      \codex{rainfall}, \codex{air\_\allowbreak{}temp\_\allowbreak{}c}, \codex{jrc\_\allowbreak{}persistent},
      \codex{lc\_\allowbreak{}built}, \codex{lc\_\allowbreak{}wetland}, \codex{lc\_\allowbreak{}open\_\allowbreak{}plus\_\allowbreak{}water},
      \codex{lc\_\allowbreak{}crop}, filtered to those present in the file, and then
      \emph{truncated to the first} \EnvChannels{}
      (\codex{dataset.py:72-73}). If fewer than \EnvChannels{} variables are
      present, the stack is padded with zero planes
      (\codex{dataset.py:74-77}).
\item From a \emph{GeoTIFF}, the first four bands are taken, zero-padded if
      the file has fewer (\codex{dataset.py:80-86}).
\item If the file cannot be read for any reason, the entire environmental
      stack is discarded and set to \codex{None}
      (\codex{dataset.py:87-88}), and every environmental channel of every
      tile becomes zero (\codex{dataset.py:249}).
\end{itemize}

\subsection{What actually happens with the real environmental NetCDF}

This is the most important fact in this section, and it is a defect in the
code rather than a property of the data.

The environmental artefact that the ingest stage produces stores the two
dynamic variables with a time dimension and the static ones without one:
\codex{rainfall} and \codex{air\_\allowbreak{}temp\_\allowbreak{}c} carry the dimensions
(\codex{time}, \codex{y}, \codex{x}), while \codex{jrc\_\allowbreak{}persistent},
\codex{lc\_\allowbreak{}built}, \codex{lc\_\allowbreak{}wetland},
\codex{lc\_\allowbreak{}open\_\allowbreak{}plus\_\allowbreak{}water} and \codex{lc\_\allowbreak{}crop} are
two-dimensional (\codex{y}, \codex{x})
(\codex{mal-core/\allowbreak{}src/\allowbreak{}mal\_\allowbreak{}core/\allowbreak{}ingest/\allowbreak{}daily\_\allowbreak{}nc.py:404-430}). The
training path reads each variable with \codex{ds[v].values.squeeze()} and
stacks the results with \codex{np.stack}
(\codex{mal-core/\allowbreak{}src/\allowbreak{}mal\_\allowbreak{}core/\allowbreak{}training/\allowbreak{}dataset.py:71-78}); \codex{squeeze}
removes unit dimensions only, so the two dynamic variables stay
three-dimensional and the two static ones stay two-dimensional, and
\codex{np.stack} then raises \codex{ValueError} because the arrays do not
all have the same shape. The \codex{except Exception} handler swallows that
error and sets the cache to \codex{None}
(\codex{mal-core/\allowbreak{}src/\allowbreak{}mal\_\allowbreak{}core/\allowbreak{}training/\allowbreak{}dataset.py:87-88}).

\textbf{The consequence is that training with the real environmental NetCDF
produces all-zero environmental channels.} The four-plane composition
described below never materialises; the network is fitted to the state and
transmission channels alone, and the four environmental input planes are
constant zeros. The prediction stage does not share the defect: its loader
takes \codex{vals[0]} when a variable is three-dimensional
(\codex{mal-core/\allowbreak{}src/\allowbreak{}mal\_\allowbreak{}core/\allowbreak{}prediction/\allowbreak{}env\_\allowbreak{}loader.py:57-62}) and
reprojects the stack onto the area-of-interest grid
(\codex{mal-core/\allowbreak{}src/\allowbreak{}mal\_\allowbreak{}core/\allowbreak{}prediction/\allowbreak{}env\_\allowbreak{}loader.py:69-84,120-146}), so
it reads real environmental planes where the training path reads zeros. The
asymmetry is silent in both directions.

Had the stack been read, the four planes that survive truncation in the
NetCDF case would be \codex{rainfall}, \codex{air\_\allowbreak{}temp\_\allowbreak{}c},
\codex{jrc\_\allowbreak{}persistent} and \codex{lc\_\allowbreak{}built} --- and the two paths
would then disagree about the fourth (Problem two).

\subsection{The grid is assumed, not checked}

The training path aligns the environmental stack to the ABM's raster by
direct cropping: it slices the cached array at the state raster's own row
and column indices, reads no transform, and performs no resampling
(\codex{mal-core/\allowbreak{}src/\allowbreak{}mal\_\allowbreak{}core/\allowbreak{}training/\allowbreak{}dataset.py:246-247}). It therefore
\emph{assumes} that the environmental file and the state raster share a
grid. That assumption is not verified anywhere, and the environmental
NetCDF carries no geographic coordinates of its own --- its \codex{x} and
\codex{y} are bare indices --- and the grid it holds is the one inherited from
the rainfall file it was assembled from
(\codex{mal-core/\allowbreak{}src/\allowbreak{}mal\_\allowbreak{}core/\allowbreak{}ingest/\allowbreak{}daily\_\allowbreak{}nc.py:196-203,404-412}), while the
state raster is on the area-of-interest grid. The training path never
compares the two, so the alignment is a property of how the artefacts were
produced rather than a checked fact. When the grids differ in
extent, resolution or origin, the crop silently samples the wrong pixels,
and any tile that extends beyond the environmental array's extent is filled
with zeros without a word
(\codex{mal-core/\allowbreak{}src/\allowbreak{}mal\_\allowbreak{}core/\allowbreak{}training/\allowbreak{}dataset.py:249}).

\subsection{Problem one: the stack does not vary in time}

There is exactly one environmental snapshot for the whole dataset, and it
serves every training pair. The real rainfall plane is a daily field, and the
environment contract lists \codex{rainfall} among the variables of the
environmental artefact (\codex{docs/\allowbreak{}concepts/\allowbreak{}generated/\allowbreak{}contracts.tex}), but the
training path has no notion of a per-day environment. Consequently the
surrogate \textbf{cannot learn the effect of a rainfall event on the next
day's state}, and a model fitted this way treats the environmental channels
as a static spatial covariate. With the real environmental NetCDF the
question is moot in the stronger sense: the channels are zero throughout
(the previous subsection), so no environmental information of any kind
reaches the network.

\subsection{Problem two: the fourth plane differs from the prediction stage's}

The prediction stage loads its environmental stack in the order
\codex{rainfall}, \codex{air\_\allowbreak{}temp\_\allowbreak{}c}, \codex{jrc\_\allowbreak{}persistent},
\codex{lc\_\allowbreak{}wetland}
(\codex{mal-core/\allowbreak{}src/\allowbreak{}mal\_\allowbreak{}core/\allowbreak{}prediction/\allowbreak{}env\_\allowbreak{}loader.py:17-18}). The training
path's fourth plane would be \codex{lc\_\allowbreak{}built}. Were the training stack read
at all, a checkpoint trained with an environmental file supplied would
receive, at inference time, a channel it was never trained on
(\codex{lc\_\allowbreak{}wetland} in place of \codex{lc\_\allowbreak{}built}), and the two
channels are not interchangeable: one is a built-surface fraction, the
other a wetland fraction. This is a defect at the training--prediction seam.
Nothing in either path detects it: the wrapper's channel-count check compares
counts, not identities (\codex{mal-core/\allowbreak{}src/\allowbreak{}mal\_\allowbreak{}core/\allowbreak{}training/\allowbreak{}wrapper.py:79-88}),
and the four planes are, as far as the code is concerned, interchangeable
numbers. As it happens, with the real environmental NetCDF the seam defect
never materialises either, because the training stack is discarded before
its composition can matter: the trained network sees four zero planes where
the prediction stage supplies four real ones.

\subsection{Problem three: the silent fallback}

If the environmental file is unreadable, the stage does not stop and does not
warn. It trains a network with no environmental context at all, and reports
nothing. The only symptom is a worse model, whose cause the operator must
infer.

With the real environmental NetCDF this fallback is not an edge case but the
normal path: the shape mismatch described above raises inside the loader,
is caught by the blanket \codex{except Exception}, and is reported to
nobody. The trained model is indistinguishable, from its artefacts, from one
trained without an environmental file at all.

% ===========================================================================
\section{The train/validation split}
\label{sec:split}

\subsection{The rule}

The split is \textbf{spatial and per-tile}. The raster's mid-column
\(c_{\mathrm{mid}} = \lfloor W/2 \rfloor\) is computed once
(\codex{mal-core/\allowbreak{}src/\allowbreak{}mal\_\allowbreak{}core/\allowbreak{}training/\allowbreak{}dataset.py:131}); a tile whose first
column index satisfies \(c \cdot P \ge c_{\mathrm{mid}}\) is a validation
tile, and every other tile is a training tile
(\codex{dataset.py:149-153}). Tiles west of the boundary train; tiles east of
it validate. The boundary is a column of the raster, so it is a geographic
statement only in the sense that the raster's own extent gives it meaning ---
it is \emph{not} a meridian, and no longitude is involved anywhere in the
code.

No size and no ratio is recorded anywhere. The split is expressed as a
column index, so the validation side receives the tile columns whose first
cell lies at or beyond the mid-column --- with \(P = \PatchSize\), roughly the
eastern half of the raster in whole tiles, but not exactly half of the cells
and not a fixed fraction. The code fixes no target proportion, and the
manifest does not record the rule that was used
(Section~\ref{sec:produces}).

\subsection{What this split does and does not guarantee}

The rule partitions space, not rollouts. The same seed contributes tiles to
both splits, and the tile immediately west of the boundary is spatially
adjacent to the tile immediately east of it. There is no held-out seed, no
buffer between the two regions, no cross-validation and no shuffling by
rollout. The validation score is therefore an \emph{in-region,
adjacent-tile} score: it measures the network's ability to generalise a few
cells across a vertical line in the same region, on the same rollout, from the
same forcing. It is not a measure of generalisation to a new area, a new
season or a new intervention.

No further defence against leakage or overfitting exists in the code. The
split rule is the only regularisation present; there is no early stopping
(Section~\ref{sec:selection}), no weight decay, no data augmentation and no
dropout.

\subsection{An empty validation split is not an error}

The trainer raises when the \emph{training} split is empty
(\codex{mal-core/\allowbreak{}src/\allowbreak{}mal\_\allowbreak{}core/\allowbreak{}training/\allowbreak{}trainer.py:47-48}) but treats an empty
validation split as legitimate: the evaluation block is guarded by a length
check, and the validation Dice stays at zero for every epoch
(\codex{trainer.py:74-82}).

The validation split \emph{can} be empty, and the mechanism is worth
stating precisely, because it is a property of the rule rather than of any
particular raster. A tile is classified by its \emph{first} column only, so
the tile that straddles the mid-column is assigned by its left edge
(\codex{dataset.py:129-131,148-151}). The validation side receives the tile
columns with \(c \ge \lceil c_{\mathrm{mid}}/P \rceil\), and the raster
contributes \(\lfloor W/P \rfloor\) tile columns; the validation split is
therefore empty whenever
\(\lceil c_{\mathrm{mid}}/P \rceil \ge \lfloor W/P \rfloor\). For a width
of \(W = 300\) cells, for instance, the tiles start at columns \(0\) and
\(128\), the mid-column is \(150\), and both tiles are training tiles, so
the validation set has no members at all.

When the validation split is empty, the validation Dice is zero at every
epoch, and the selection rule is a tie at every epoch: the comparison is
\codex{>=} and the running best starts at \(0.0\)
(\codex{trainer.py:57,86-88}), so \codex{0.0 >= 0.0} holds every time.
The checkpoint is therefore overwritten at \emph{every} epoch and
\codex{best\_model.pt} ends up holding the \textbf{last} epoch's weights,
not the first one's. The training loss of that last epoch is not necessarily
the best one observed.

% ===========================================================================
\section{The architecture and the loss}
\label{sec:architecture}

\subsection{Architecture}

The network is a standard U-Net with four resolution levels
(\codex{mal-core/\allowbreak{}src/\allowbreak{}mal\_\allowbreak{}core/\allowbreak{}training/\allowbreak{}model.py:28-46}).

\begin{itemize}
\item \textbf{Encoder.} Four stages, with the widths recorded in the
      generated parameter table below. Each
      stage is a \emph{double convolution block}: a \(3\times3\) convolution
      without bias, batch normalisation, a ReLU, and then the same three
      operations again (\codex{model.py:14-25}). A \(2\times2\) max-pool
      halves the resolution between stages (\codex{model.py:37}).
\item \textbf{Bottleneck.} A double convolution block of double the last
      encoder width, at the coarsest resolution
      (\codex{model.py:38}).
\item \textbf{Decoder.} For each width in reverse order, a transposed
      convolution with stride two upsamples, the corresponding encoder
      feature is concatenated onto it, and a double convolution block
      processes the result (\codex{model.py:39-44}).
\item \textbf{Head.} A \(1\times1\) convolution from the first width to the
      number of output channels (\codex{model.py:46}). There is no activation
      on the head: the output is an unbounded real tensor.
\end{itemize}

The generated table records the constructor's declared shape
parameters:

% Generated by tools/gen_concepts_values.py -- do not edit.
\paragraph{The surrogate's architecture}
\begin{longtable}{@{}ll@{}}
\toprule
Element & Value \\
\midrule
\endhead
\code{in\_channels} & 6 \\
\code{out\_channels} & 2 \\
\code{channels} & (32, 64, 128, 256) \\
\bottomrule
\end{longtable}

\paragraph{The training loop's defaults}
\begin{longtable}{@{}ll@{}}
\toprule
Argument & Default \\
\midrule
\endhead
\code{epochs} & 50 \\
\code{batch\_size} & 16 \\
\code{lr} & 0.001 \\
\code{device} & None \\
\code{subsample} & 1.0 \\
\code{preload} & False \\
\code{include\_transmission} & True \\
\code{env\_path} & None \\
\bottomrule
\end{longtable}


If the encoder's resolution does not match the skip feature's, the mismatch
is repaired by bilinear interpolation to the skip's size
(\codex{model.py:55-57}). With \(P = \PatchSize\) and four pooling stages the
two agree, so the interpolation branch is not exercised during training. In
the inference wrapper a different repair is used: the input is padded with
replicate padding to a multiple of \PoolMult{}
(\codex{mal-core/\allowbreak{}src/\allowbreak{}mal\_\allowbreak{}core/\allowbreak{}training/\allowbreak{}wrapper.py:90-98}) and the output is
cropped back (\codex{wrapper.py:102-104}). \PoolMult{} is the product of the
four pooling factors.

\subsection{The loss}

\begin{equation}
 \mathcal{L} \;=\; \mathrm{MSE}\!\left(\hat{y}, y\right)
 \;+\; \DiceWeight \cdot \mathrm{Dice}_{\mathrm{soft}}\!\left(\hat{y}_{0}, y_{0}\right),
 \label{eq:loss}
\end{equation}

where the mean squared error is taken over \emph{all} channels
(\codex{model.py:79}) and the soft-Dice term is computed on \emph{channel 0
only} (\codex{model.py:80}) --- the first state band, which under contract
\StateContract{} holds raw adult counts and is the sparse one. The rationale
is stated in the code itself: the Dice term focuses on the sparse channel.

The soft-Dice term smooths both arguments through a logistic function before
forming the ratio (\codex{model.py:69-73}):

\begin{equation}
 p = \sigma\!\left(\frac{\hat{y}-\theta}{\tau}\right),\qquad
 t = \sigma\!\left(\frac{y-\theta}{\tau}\right),\qquad
 \mathrm{Dice}_{\mathrm{soft}} = \frac{2\sum pt + \varepsilon}{\sum p + \sum t + \varepsilon},
\end{equation}

with temperature \(\tau = \DiceTau\), threshold \(\theta = \DiceThresh\), and
stabiliser \(\varepsilon = \DiceEps\). The sum runs over the whole batch and
the whole tile, not over pixels individually: a tile is scored as one object.

\subsection{Units of the threshold --- and why they matter}

Under state contract \StateContract{} the state bands hold raw adult counts,
so a threshold of \(\DiceThresh\) adult counts means
\emph{one adult or more in the cell} (the code says so in as many words at
\codex{model.py:66-67}). This is a deliberate and consequential choice: the
Dice term is effectively an occupancy term --- ``does this cell contain
mosquitoes tomorrow'' --- while the MSE term, whose gradient contribution
grows with the magnitude of the adult counts, is an exactness term. Which of
the two terms dominates the gradient in practice depends on the count
magnitudes in the training rasters, which this document does not fix; the
code does not record the balance, and no claim about it is made here. A
retrained checkpoint from before the
retirement of the capacity normalisation would have a different meaning for
the same number, which is one reason the manifest's silence about the
contract version (Section~\ref{sec:produces}) is a real omission rather than
a cosmetic one.

\subsection{The validation metric}

The validation metric is a \emph{hard} Dice on the same channel, using a
strict comparison against the same threshold
(\codex{mal-core/\allowbreak{}src/\allowbreak{}mal\_\allowbreak{}core/\allowbreak{}training/\allowbreak{}model.py:86-95}), averaged over the
validation batches (\codex{mal-core/\allowbreak{}src/\allowbreak{}mal\_\allowbreak{}core/\allowbreak{}training/\allowbreak{}trainer.py:75-82}).

It is important to be clear about what this number characterises. It is
computed on \emph{one band of one output}: the first state channel. It says
nothing about the second state band (\codex{adult\_\allowbreak{}cell\_\allowbreak{}count}), which is
trained only through the MSE term, and nothing about the four transmission
bands, which are also trained only through the MSE term. Their accuracy is
never measured anywhere. The reported ``best validation Dice'' is therefore a
statement about occupancy of the first state band, not about the surrogate.

% ===========================================================================
\section{Optimisation and model selection}

\label{sec:selection}

\subsection{Optimiser}

The optimiser is Adam (\codex{mal-core/\allowbreak{}src/\allowbreak{}mal\_\allowbreak{}core/\allowbreak{}training/\allowbreak{}trainer.py:55})
with the learning rate given by the caller; the packaged default is the one
in the generated table above. There is no learning-rate schedule, no weight
decay, no momentum configuration beyond Adam's own defaults, no gradient
clipping and no warm-up.

\subsection{What stands in for early stopping}

Nothing. The epoch loop has exactly two conditionals, and both concern
checkpoint writing (\codex{trainer.py:86-91}). There is no patience counter,
no early-stopping criterion, no minimum delta, and the loop always runs the
full number of epochs. The only mechanism that resembles early stopping is
the running maximum of the validation Dice, described next.

\subsection{The selection rule}

\begin{quote}
The checkpoint is rewritten whenever the validation Dice is \emph{at least}
the best value seen so far, and unconditionally at the first epoch
(\codex{mal-core/\allowbreak{}src/\allowbreak{}mal\_\allowbreak{}core/\allowbreak{}training/\allowbreak{}trainer.py:86-88}).
\end{quote}

Three points follow from the exact form of that comparison.

\begin{enumerate}
\item The comparison is \codex{>=}, so a \emph{tie} also overwrites the
      checkpoint. A later epoch that merely equals the best replaces it.
\item The write at epoch zero is unconditional, so a checkpoint always exists
      even if every subsequent validation score is worse.
\item Because the loop never stops, ``best'' is a running maximum over all
      epochs. The final \codex{best\_\allowbreak{}model.pt} is the weights of the last
      epoch that attained the maximum, not an early-stopped model, and the
      training loss of that epoch may be far from the best loss observed.
\end{enumerate}

There is a fourth consequence, and it is the one that bites in practice. The
running best is initialised to \(0.0\)
(\codex{mal-core/\allowbreak{}src/\allowbreak{}mal\_\allowbreak{}core/\allowbreak{}training/\allowbreak{}trainer.py:57}), and the
comparison is \codex{>=}, so a validation Dice of exactly zero is a tie with
the running best at \emph{every} epoch: \codex{0.0 >= 0.0} holds. The
checkpoint is therefore rewritten on every epoch, and
\codex{best\_\allowbreak{}model.pt} ends up holding the \emph{last} epoch's
weights rather than the first one's. A zero validation Dice is what an empty
validation split produces (Section~\ref{sec:selection}), so a run whose
validation set is empty selects its final epoch and calls it ``best'' without
saying so (Section~\ref{sec:split}).

\subsection{What the operator sees}

One line per epoch, carrying the mean training loss and the validation Dice
(\codex{mal-core/\allowbreak{}src/\allowbreak{}mal\_\allowbreak{}core/\allowbreak{}training/\allowbreak{}trainer.py:84}); and a final line
carrying the best validation Dice, which is also the function's return value
(\codex{trainer.py:108-109}), echoed by the command-line wrapper
(\codex{mal-core/\allowbreak{}src/\allowbreak{}mal\_\allowbreak{}core/\allowbreak{}cli.py:806}).

\subsection{Reproducibility}

Training is \emph{not} reproducible end to end, and the reason is worth
stating precisely rather than as a blanket complaint. Only the tile
subsample is seeded (\codex{mal-core/\allowbreak{}src/\allowbreak{}mal\_\allowbreak{}core/\allowbreak{}training/\allowbreak{}dataset.py:167}).
The global torch random number generator is never seeded, so the weight
initialisation differs between runs; the training loader shuffles with an
unseeded generator (\codex{dataset.py:281-283}); and no seed, RNG state or
epoch counter is stored in the checkpoint. Two runs over the same data will
produce different weights and, almost surely, different validation Dice
values. Only the \emph{set of tiles} is reproducible.

\subsection{Device}

The device is resolved as CUDA if available, else Apple MPS, else CPU
(\codex{mal-core/\allowbreak{}src/\allowbreak{}mal\_\allowbreak{}core/\allowbreak{}training/\allowbreak{}trainer.py:29-35}); the inference
wrapper repeats the same resolution (\codex{wrapper.py:29-37}). The choice is
not recorded in any artefact, so a checkpoint carries no hint of the hardware
it was fitted on.

% ===========================================================================
\section{Checkpoints, and how one is loaded}

\subsection{What is inside a checkpoint}

All three checkpoint files are written with the same call: a bare
\codex{torch.save} of the network's \codex{state\_\allowbreak{}dict}
(\codex{mal-core/\allowbreak{}src/\allowbreak{}mal\_\allowbreak{}core/\allowbreak{}training/\allowbreak{}trainer.py:88,91,93}). A
\codex{state\_\allowbreak{}dict} holds parameters and buffers only. It does \emph{not} hold:

\begin{itemize}
\item the optimiser state, so training cannot be resumed;
\item the epoch number or the metric history;
\item the random number generator's state;
\item the architecture description --- the channel counts and the widths must
      be recovered from elsewhere;
\item the loss configuration, the split rule, or any data provenance;
\item the ABM state-contract version of the rasters it was fitted to.
\end{itemize}

A checkpoint is therefore an inference artefact and nothing else. Resuming
training from one is impossible.

\subsection{How a checkpoint is loaded}

The wrapper's constructor performs the load eagerly, in four steps
(\codex{mal-core/\allowbreak{}src/\allowbreak{}mal\_\allowbreak{}core/\allowbreak{}training/\allowbreak{}wrapper.py:39-66}):

\begin{enumerate}
\item \textbf{The manifest first.} If a \codex{model.yaml} sits beside the
      checkpoint, its \codex{in\_\allowbreak{}channels} and \codex{out\_\allowbreak{}channels} are read
      and, where the caller did not supply them, used
      (\codex{wrapper.py:39-46}).
\item \textbf{Shape inference as a fallback.} If the channel counts are still
      unknown, they are inferred from the tensors themselves: the first
      encoder convolution's weight tensor carries the input width, and the
      head's weight tensor carries the output width
      (\codex{wrapper.py:51-59}). Fixed fallbacks of
      \codeconst{training.fallback_in_channels} and
      \codeconst{training.fallback_out_channels} apply if even
      those tensors are absent.
\item \textbf{Rebuild and load.} The network is constructed with the
      resolved counts and the weights are loaded with a strict
      \codex{load\_\allowbreak{}state\_\allowbreak{}dict}, then the network is put into evaluation mode
      (\codex{wrapper.py:64-67}). A missing or mismatched checkpoint raises at
      construction --- there is no lazy path and no partially initialised
      model.
\item \textbf{Input adaptation.} At call time, if the channels presented sum
      to fewer than the model's input width, the input is zero-padded; if
      they sum to more, the input is truncated
      (\codex{wrapper.py:79-88}).
\end{enumerate}

Step 4 deserves emphasis: the wrapper \emph{silently adapts} to a channel
mismatch rather than refusing. A network trained without transmission, fed
the six-channel state stack that the prediction stage assembles, will receive
its first six channels and ignore the rest; a network trained with
transmission but fed only a two-channel state will receive four zero planes
it was never trained to interpret. Neither case produces a warning, a log
line, or an exception. The prediction stage does pass the manifest's counts
to the wrapper (\codex{mal-core/\allowbreak{}src/\allowbreak{}mal\_\allowbreak{}core/\allowbreak{}prediction/\allowbreak{}registry.py:110-114}),
so in the normal path the widths agree; the adaptation exists for the
mismatch cases, and it hides them.

% ===========================================================================
\section{The seam with the prediction stage}

\subsection{How the prediction stage finds the model}

The registry scans two roots --- the models directory and the runs directory
--- for files named \codex{model.yaml}, and keys each model as
\codex{name@version}
(\codex{mal-core/\allowbreak{}src/\allowbreak{}mal\_\allowbreak{}core/\allowbreak{}prediction/\allowbreak{}registry.py:54-74}). Loading a model
by that key builds the wrapper from the manifest's declared checkpoint
(\codex{registry.py:96-118}). Finding a model is not the same as using it. The prediction command's
\codex{--model} option defaults to \codex{"dummy"}
(\codex{mal-core/\allowbreak{}src/\allowbreak{}mal\_\allowbreak{}core/\allowbreak{}cli.py:112}), and the registry answers
that name with the dummy without consulting the scan at all
(\codex{mal-core/\allowbreak{}src/\allowbreak{}mal\_\allowbreak{}core/\allowbreak{}prediction/\allowbreak{}registry.py:96-97}); a trained checkpoint is
therefore applied only when the operator names its registry key explicitly.
If no trained model exists, the registry
returns a dummy model that emits zeros of the input's width
(\codex{registry.py:41-46,96-97,115-116}); a prediction made with it is a raster of
zeros, and the sidecar still names a model, so the reader of the output must
look at the model name to tell.

\subsection{What the prediction stage hands to the surrogate}

The prediction stage assembles two tensors and calls the surrogate exactly
once (\codex{mal-core/\allowbreak{}src/\allowbreak{}mal\_\allowbreak{}core/\allowbreak{}prediction/\allowbreak{}predictor.py:84-91}):

\begin{itemize}
\item \textbf{The state stack.} Loaded by the state loader with transmission
      requested (\codex{predictor.py:54-56}), giving two vector bands and,
      when the transmission raster is found, four transmission bands. If
      nothing is found the stage substitutes zeros of six channels
      (\codex{predictor.py:55-56}).
\item \textbf{The environmental stack.} Four planes on the area-of-interest
      grid (\codex{predictor.py:58-61}), falling back to zeros if the
      environment cannot be loaded (\codex{predictor.py:60-61}).
\end{itemize}

The answer is then aggregated to the requested scale --- no aggregation for
the regional scale, a mean-pool for the national scale, a single mean for the
continental scale (\codex{mal-core/\allowbreak{}src/\allowbreak{}mal\_\allowbreak{}core/\allowbreak{}prediction/\allowbreak{}aggregator.py:51-71})
--- and written as a risk GeoTIFF with band descriptions and a JSON sidecar
(\codex{predictor.py:94-150}). With six output channels the band names are the
predicted adult patch count, the predicted adult cell count, the predicted
human prevalence, incidence, infectious pressure and transmission focus
(\codex{predictor.py:30-41}).

\subsection{Four mismatches across the seam}

The training and inference paths do not agree on four points. All four are
visible in the code; none is recorded anywhere in the repository as a known
issue.

\begin{enumerate}
\item \textbf{The environmental plane order differs}, as established in
      Section~\ref{sec:env}: \codex{lc\_\allowbreak{}built} at training,
      \codex{lc\_\allowbreak{}wetland} at inference.
\item \textbf{The environmental grids differ.} At training the environmental
      planes are those of the environmental NetCDF, whose grid is inherited
      from the rainfall file and which carries no geographic coordinates
      (\codex{mal-core/\allowbreak{}src/\allowbreak{}mal\_\allowbreak{}core/\allowbreak{}ingest/\allowbreak{}daily\_\allowbreak{}nc.py:196-203,404-412}); at
      inference the same variables are reprojected onto the
      area-of-interest grid
      (\codex{mal-core/\allowbreak{}src/\allowbreak{}mal\_\allowbreak{}core/\allowbreak{}prediction/\allowbreak{}env\_\allowbreak{}loader.py:69-84,120-146}). The two
      sides therefore do not see the same environmental quantities, and the
      training side does not resample anything at all
      (\codex{mal-core/\allowbreak{}src/\allowbreak{}mal\_\allowbreak{}core/\allowbreak{}training/\allowbreak{}dataset.py:246-247}).
\item \textbf{The state stack is resampled at inference and cropped at
      training.} The state loader reprojects the ABM's raster onto the
      area-of-interest grid with bilinear interpolation
      (\codex{mal-core/\allowbreak{}src/\allowbreak{}mal\_\allowbreak{}core/\allowbreak{}prediction/\allowbreak{}state\_\allowbreak{}loader.py:99-107,140-148}), while
      the training path crops raw ABM pixels and never interpolates
      (\codex{mal-core/\allowbreak{}src/\allowbreak{}mal\_\allowbreak{}core/\allowbreak{}training/\allowbreak{}dataset.py:210-214}). The
      surrogate is therefore applied, at prediction time, to an interpolated
      grid whose local statistics differ from the crisp pixels it was fitted
      on --- and an adult count, which is an integer number of individuals,
      is being bilinearly blended.
\item \textbf{The channel counts agree only under a specific history.} The
      prediction stage requests the transmission bands as well
      (\codex{predictor.py:54-56}), so the stack it assembles has six
      channels when the transmission raster is found and two when it is not
      (\codex{predictor.py:55-56}). A surrogate trained with transmission
      matches the first case and only that case; in the second it is fed a
      two-channel stack and padded to the width it expects with zeros
      (\codex{wrapper.py:79-88}). A checkpoint trained without transmission
      is silently adapted by the same rule, in the opposite direction, by
      truncation.
\end{enumerate}

A fifth difference belongs to the same seam, though it is a consequence of
the first rather than an independent mismatch: with the real environmental
NetCDF the training side's four environmental channels are all zeros
(Section~\ref{sec:env}), while the prediction side's four are real
environmental planes. A surrogate fitted on the training side has never seen
a non-zero environmental channel at all.

\subsection{The surrogate is applied once}

There is no rollout loop. The prediction stage calls \codex{model.predict}
once and aggregates the result
(\codex{mal-core/\allowbreak{}src/\allowbreak{}mal\_\allowbreak{}core/\allowbreak{}prediction/\allowbreak{}predictor.py:89-92}). The output is a
\emph{risk map}, not a projection: it is one local transition applied to one
snapshot, and it does not yield trajectories, prevalence over time, or
intervention scenarios unless the rollouts the surrogate learned from
contained them.

% ===========================================================================
\section{Declarations of absence}

\label{sec:absences}

This section records what the stage does \emph{not} do, and does \emph{not}
have. Each item was verified against the tree, not inferred.

\subsection{There is no calibration or acceptance harness}

The pipeline has a scoring framework --- a directory of scorers, a registry
of their weights, a composite score and a scorecard --- under
\codex{mal-core/\allowbreak{}src/\allowbreak{}mal\_\allowbreak{}core/\allowbreak{}scoring/\allowbreak{}} (with the
scorers themselves in \codex{mal-core/\allowbreak{}src/\allowbreak{}mal\_\allowbreak{}core/\allowbreak{}scoring/\allowbreak{}scorers/\allowbreak{}}).
The training stage has
no counterpart. Searching the training package
(\codex{mal-core/\allowbreak{}src/\allowbreak{}mal\_\allowbreak{}core/\allowbreak{}training/\allowbreak{}}) and its test file
(\codex{mal-core/\allowbreak{}tests/\allowbreak{}test\_\allowbreak{}training\_\allowbreak{}pipeline.py}) finds no scorer, no
thresholds file, no composite, no tolerance band, no scorecard, and no
acceptance threshold of any kind. The only notion of ``best'' in the stage is
the running maximum of the validation Dice described in
Section~\ref{sec:selection}.

The test file that does exist checks plumbing rather than fidelity: the
network's output shapes for the two channel configurations, that the loss
returns a non-negative tensor, that the dataset yields the expected tensor
shapes, and that an end-to-end training run produces a checkpoint and a
manifest that the wrapper can load
(\codex{mal-core/\allowbreak{}tests/\allowbreak{}test\_\allowbreak{}training\_\allowbreak{}pipeline.py:38-114}).

Two artefacts outside the package are worth recording here, because they are
easy to mistake for the harness that does not exist.

\begin{itemize}
\item \codex{mal-execution/\allowbreak{}scripts/\allowbreak{}validate\_\allowbreak{}unet.py} contains a hard-coded
      acceptance threshold on the validation Dice (a Dice of 0.6), prints a
      pass/fail verdict and writes it to a text file
      (\codex{mal-execution/\allowbreak{}scripts/\allowbreak{}validate\_\allowbreak{}unet.py:47-57}). It is a
      standalone script, not reachable from the command line, not part of the
      package, and not a gate that anything consults. It also constructs the
      network with the constructor's default channel counts
      (\codex{validate\_\allowbreak{}unet.py:31}), so it cannot load a transmission-enabled
      checkpoint.
\item \codex{mal-execution/\allowbreak{}scripts/\allowbreak{}train\_\allowbreak{}unet\_\allowbreak{}subsample.py} imports the
      network from a module path that does not exist in the package
      (\codex{train\_\allowbreak{}unet\_\allowbreak{}subsample.py:11}); the import fails before any
      training happens. It also pairs a day with the day seven days later
      (\codex{train\_\allowbreak{}unet\_\allowbreak{}subsample.py:65-66}), which is a different learning
      problem from the one the packaged trainer solves.
\end{itemize}

The third script, \codex{mal-execution/\allowbreak{}scripts/\allowbreak{}train\_\allowbreak{}unet.py}, calls the
packaged trainer through the public API and is consistent with it.

\subsection{The surrogate is never validated against the ABM}

No comparison between a surrogate's output and a subsequent ABM step exists
anywhere in the tree; no accuracy floor, no regression test that a retrained
checkpoint stays within a tolerance of a previous one. The surrogate's output
goes straight into the risk map
(\codex{mal-core/\allowbreak{}src/\allowbreak{}mal\_\allowbreak{}core/\allowbreak{}prediction/\allowbreak{}predictor.py:89}). Fidelity is
therefore \emph{ungated}: nothing in the pipeline would notice a checkpoint
that had learned the wrong thing.

\subsection{Other absences, itemised}

\begin{itemize}
\item No learning-rate schedule, no early stopping, no patience rule, no
      gradient clipping (\codex{mal-core/\allowbreak{}src/\allowbreak{}mal\_\allowbreak{}core/\allowbreak{}training/\allowbreak{}trainer.py:57-93}).
\item No seeding of the global torch generator and no shuffling seed
      (\codex{dataset.py:167,281-283}).
\item No resumption from a checkpoint: no optimiser state, no epoch counter,
      no RNG state in the checkpoint (\codex{trainer.py:88,91,93}).
\item No data provenance in the manifest beyond a free-text description
      (\codex{trainer.py:96-105}).
\item No record of the state-contract version in the manifest
      (\codex{trainer.py:96-105}).
\item No record of the device, the seed, the split rule or the patch size in
      any artefact.
\item No environmental channel-identity check between training and inference
      (\codex{dataset.py:72-73} against
      \codex{mal-core/\allowbreak{}src/\allowbreak{}mal\_\allowbreak{}core/\allowbreak{}prediction/\allowbreak{}env\_\allowbreak{}loader.py:17-18}).
\item No no-data mask: no-data becomes zero, indistinguishable from an empty
      cell (\codex{dataset.py:181-182}).
\item No per-day environmental stack, so no rainfall-event response can be
      learned (\codex{dataset.py:59,66-88}).
\item No handling of the mixed dimensionality of the real environmental
      NetCDF: the stack is discarded wholesale by a blanket exception, so the
      environmental channels of a real training run are all zero
      (\codex{dataset.py:71-78,87-88}).
\item No grid check between the environmental file and the state raster: the
      crop is positional, and tiles beyond the environmental extent are
      silently zeroed (\codex{dataset.py:246-249}).
\item No remainder cells: a raster dimension not divisible by \(P\) loses its
      edge strip from training (\codex{dataset.py:129-130}).
\item No held-out seed, no buffered split, no cross-validation
      (\codex{dataset.py:149-153}).
\end{itemize}

\subsection{What is not answered by the code}

The following are questions the code does not settle. They are recorded here
rather than answered by guesswork.

\begin{enumerate}
\item Whether the mid-column split is \emph{intended} to be a regional
      in-sample score, or whether a spatially buffered split was expected,
      is not recorded anywhere.
\item Whether the absence of training-data provenance in the manifest is
      deliberate or an oversight is not recorded.
\item Whether an early-stopping rule, a learning-rate schedule or a patience
      counter is wanted, or their absence is a decision, is not recorded.
\item Whether end-to-end reproducibility is a requirement is not recorded.
\item Whether the manifest should carry the ABM state-contract version is not
      recorded.
\item Whether the environmental channel-order mismatch, the discarded
      environmental stack and the inference-time bilinear reprojection are
      known defects is not stated anywhere in the
      repository.
\item Whether the surrogate's fidelity is intended to be gated against the
      ABM at some point, or whether the stage is knowingly ungated, is not
      recorded.
\item Whether the training hyperparameters are intended to move into the
      shared configuration file is a decision that has not been taken.
\item No written prohibition exists on the uses the surrogate must not be put
      to; the restrictions in Section~\ref{sec:purpose} are inferences from
      the code's one-step, one-call structure.
\end{enumerate}

% ===========================================================================
\section{Evidence appendix}

\subsection{How the values in this document are sourced}

Three sources, in increasing order of authority for this stage:

\begin{enumerate}
\item \textbf{The shared configuration.} Values that exist in
      \codex{mal-core/\allowbreak{}config/\allowbreak{}model\_\allowbreak{}constants.yaml} are referenced through the
      \code{\textbackslash mskey} macro, so that a renamed key makes this
      document fail to build rather than silently print a stale number. The
      values of that kind used here are the state contract version
      (\codex{contracts.cogs\_\allowbreak{}state}), the environment contract version
      (\codex{contracts.env\_\allowbreak{}nc}) and the no-data sentinel
      (\mskey{ingest.nodata_sentinel}).
\item \textbf{The generated tables.} The surrogate's declared constructor
      parameters and the training loop's defaults are emitted by
      \codex{tools/\allowbreak{}gen\_\allowbreak{}concepts\_\allowbreak{}values.py} --- the training table by
      \codex{emit\_\allowbreak{}training\_\allowbreak{}params} at line 463, which reads the defaults with
      \codex{\_\_function\_\allowbreak{}defaults} at lines 423--460. That helper parses the
      source \emph{statically}, with Python's \codex{ast} module and without
      importing the training package (its own docstring says so), so the
      values are read from the code text rather than by introspection of
      imported objects. The tables are included above via
      \codex{\textbackslash input} of
      \codex{../\allowbreak{}generated/\allowbreak{}params-training.tex}. The artefact contracts are
      emitted from \codex{docs/\allowbreak{}concepts/\allowbreak{}contracts.yaml} and included via
      \codex{../\allowbreak{}generated/\allowbreak{}contracts.tex}.
\item \textbf{The code literals.} Everything else in this document is a
      literal in the source of the training package. The shared configuration
      file contains no key that mentions the training stage, the surrogate,
      the network, the loss or the number of epochs --- this was checked by
      searching that file --- so these numbers have no single source of truth
      in the shared configuration file. They are read out of the code text by
      \codex{tools/\allowbreak{}gen\_\allowbreak{}concepts\_\allowbreak{}values.py} and emitted as
      \codex{\textbackslash codeconst} macros in
      \codex{docs/\allowbreak{}concepts/\allowbreak{}generated/\allowbreak{}code-macros.tex}, each carrying the
      file and line it was read from; the macros declared at the top of this
      source file are short aliases for them, so an unknown key makes this
      document fail to build rather than print a stale number.
\end{enumerate}

\subsection{Code literals used in this document}

Each of these is a reading of the code, verified at the cited line.

\begin{center}
\begin{longtable}{@{}lp{2.4cm}lp{3.0cm}p{2.6cm}p{4.0cm}@{}}
\toprule
Quantity & Value & Where it is written & What it controls \\
\midrule
\endhead
Tile side \(P\) & \PatchSize & \codex{dataset.py:19} & The side of a training tile in raster cells \\
Vector channels & \VectorChannels & \codex{dataset.py:20} & Declared; unused in practice \\
Transmission channels & \TransChannels & \codex{dataset.py:21} & Declared; unused in practice \\
Environmental channels & \EnvChannels & \codex{dataset.py:22} & The width of the environmental stack \\
Dice weight & \DiceWeight & \codex{model.py:81} & The weight of the soft-Dice term in the loss \\
Dice temperature \(\tau\) & \DiceTau & \codex{model.py:63} & The smoothing of the soft-Dice term \\
Dice threshold \(\theta\) & \DiceThresh & \codex{model.py:63,86} & The count threshold in the soft and hard Dice \\
Dice stabiliser \(\varepsilon\) & \DiceEps & \codex{model.py:73,95} & The numerical floor in both Dice ratios \\
Subsample seed & \SubSeed & \codex{dataset.py:167} & The reproducibility of the tile subsample \\
Checkpoint period & \CkptPeriod & \codex{trainer.py:90-91} & How often a periodic checkpoint is written \\
Padding multiple & \PoolMult & \codex{wrapper.py:93-94} & The size to which the inference input is padded \\
\bottomrule
\end{longtable}
\end{center}

\subsection{Implementation notes a replicator needs}

These are behaviours invisible from the description above, and they are the
ones that will bite a reimplementation.

\begin{enumerate}
\item \textbf{Adjacency is in the sorted list of present days, not in
      integer days.} A run directory with days \(1,2,9,10\) produces the
      pair \((2,9)\), and the network is being taught a one-day transition
      from a seven-day gap (\codex{dataset.py:135-146}).
\item \textbf{Bands are read positionally.} The training path never checks a
      band description against an expected name
      (\codex{dataset.py:178-179}); a reordering of the ABM's bands would
      silently permute the target.
\item \textbf{The seed fallback is one.} A file whose name carries neither
      \codex{seedN} nor \codex{\_\allowbreak{}dayN} is filed under seed one, day zero
      (\codex{dataset.py:27-34}).
\item \textbf{The subsample is drawn per split.} The training and validation
      datasets each seed the generator to the same value and each draw from
      their own list (\codex{dataset.py:166-170}), so the two subsets are
      reproducible but different.
\item \textbf{An empty validation split is silently tolerated} and yields a
      validation Dice of zero for every epoch; because the running best
      starts at \(0.0\) and the comparison is \codex{>=}, the checkpoint is
      rewritten every epoch and the \emph{last} epoch's weights end up in
      \codex{best\_\allowbreak{}model.pt}
      (\codex{trainer.py:57,74-88}).
\item \textbf{Ties overwrite the best checkpoint}, because the comparison is
      \codex{>=} and not \codex{>} (\codex{trainer.py:86}).
\item \textbf{The environment is read once and cached for the whole
      dataset}, and a read failure replaces it with nothing at all
      (\codex{dataset.py:66-88}).
\item \textbf{The real environmental NetCDF trips that fallback.} Its two
      dynamic variables are three-dimensional and its statics are not, so
      \codex{np.stack} raises and the blanket \codex{except Exception}
      discards the whole stack; the four environmental channels of every
      training tile are consequently zero
      (\codex{dataset.py:71-78,87-88}; the dimensions are set in
      \codex{mal-core/src/mal\_core/ingest/daily\_nc.py:404-430}).
\item \textbf{The environmental crop is positional}, so the training path
      assumes the environmental file and the state raster share a grid and
      samples the wrong pixels whenever they do not; tiles beyond the
      environmental extent are zeroed without comment
      (\codex{dataset.py:246-249}).
\item \textbf{The inference wrapper adapts to channel mismatches by padding
      with zeros or truncating}, silently (\codex{wrapper.py:79-88}).
\item \textbf{The inference wrapper pads to a multiple of \PoolMult{} with
      replicate padding} and crops the result back
      (\codex{wrapper.py:90-104}), so the border cells of the output are
      computed from replicated values rather than from neighbouring data.
\item \textbf{The dummy model emits zeros of the input's width}
      (\codex{registry.py:41-46}); a prediction produced without a trained
      model is a raster of zeros with a sidecar that still names a model.
\item \textbf{A trained checkpoint is never selected automatically.} The
      prediction command's \codex{--model} default is \codex{"dummy"}
      (\codex{mal-core/src/mal\_core/cli.py:112}) and the registry answers
      that name directly (\codex{registry.py:96-97}), so the surrogate runs
      only when its registry key is passed explicitly.
\end{enumerate}

\subsection{What is verified, and what is not}

\textbf{Verified by reading the code in this repository, at the cited lines:}
the file discovery and naming conventions; the seed and day parsing; the
pairing rule; the tile grid and the dropped remainder; the mid-column split
and its exact inequality, including the empty-validation degenerate case;
the environmental variable order and truncation, and the fact that the real
environmental NetCDF is discarded by the shape mismatch; the silent
environment fallback; the positional environmental crop and the
zero-fill beyond the environmental extent; the no-data-to-zero conversion;
the
subsample's determinism; the architecture and the loss composition; the
threshold's unit under state contract \StateContract; the optimiser, the
absence of a schedule and of early stopping; the selection rule's exact
comparison, including the \codex{>=} tie-overwrite and its consequence for
an all-zero validation Dice; the checkpoint contents; the manifest's keys and
omissions; the
loading order and the shape-inference fallback; the silent input adaptation;
the registry's discovery and the dummy model's behaviour; the prediction
stage's single call and its aggregation; and the absence of any calibration
harness, scorer, threshold or accuracy gate in the training package and its
tests.

\textbf{Not verified, and not claimed:} the numerical fidelity of any
checkpoint (no checkpoint and no trained model were exercised here, and no
accuracy comparison against the ABM exists in the tree to compare against);
the behaviour of the stage on real ABM output, as opposed to the synthetic
rasters used by the test file; the correctness of the ABM's own
output contract, which this document takes as given from the ABM's
description; and the intent behind the absences listed in
Section~\ref{sec:absences}, which the code does not record.

Three further statements are statistical expectations rather than code
facts, and are marked as such where they appear: \emph{what a fitted
network generalises to} beyond the rollouts it saw, since the code records
no such measurement; \emph{the relative gradient weight of the MSE and
Dice terms}, which depends on count magnitudes the code does not fix and
records no trace of beyond the printed training loss; and \emph{the fidelity
of the surrogate}, which is ungated by construction.

\textbf{The retired specification is not used as a source.} The stage's
former specification contains several statements that the code contradicts:
it describes a fixed seven-day step where the code pairs consecutive days;
it describes a narrow file-name glob where the code globs loosely; it gives
different values for the Dice temperature, the Dice threshold and the
selection comparison; and it presents the environmental zero-fill as a
standing limitation where the code treats it as a fallback. None of its
claims is cited as authority here; the code is.

\end{document}
